\documentclass[ecta,nameyear,draft]{econsocart}
\usepackage{amsmath,amsfonts,amssymb,amsthm,booktabs,array,longtable}
\usepackage[T1]{fontenc}
\usepackage{mathpazo,microtype,needspace,xurl}
\RequirePackage[colorlinks,citecolor=blue,linkcolor=blue,urlcolor=blue]{hyperref}
\startlocaldefs
\newtheorem{theorem}{Theorem}
\newtheorem{proposition}[theorem]{Proposition}
\newtheorem{lemma}[theorem]{Lemma}
\newtheorem{corollary}[theorem]{Corollary}
\newcommand{\E}{\mathbb E}
\newcommand{\R}{\mathbb R}
\endlocaldefs
\setlength{\emergencystretch}{3em}
\makeatletter\def\copyright@text{Prepared for referee review}\makeatother

% Automatically derived from verified result records.
\providecommand{\NBOServices}{96}
\providecommand{\NBOAttained}{80}
\providecommand{\NBOExhausted}{16}
\providecommand{\NBOFixedActors}{5}
\providecommand{\NBOFixedAliases}{32}
\providecommand{\NBOFactorialAnswers}{3160}
\providecommand{\NBOCenteredGap}{0.03979}
\providecommand{\NBOExtraChanges}{12}

\setcounter{secnumdepth}{2}
\usepackage{xr-hyper}
\externaldocument[M-]{build/development}[development.pdf]
\begin{document}\begin{frontmatter}
\title{Supplement to ``Neural Bellman Operators''}
\runtitle{Neural Bellman Operators: Supplement}
\begin{aug}\author[id=au1,addressref={add1}]{\fnms{Qian}~\snm{QI}}\address[id=add1]{Peking University}\end{aug}
\begin{abstract}
This supplement gives the primitive-domain verification, numerical and statistical details, complete resource records and direct economic comparisons for the retained R49 evidence and the R54 revision. The reference-advantage, terminal-improvement and finite-sweep accuracy proofs are included here; all earlier proofs remain in the complete development supplement. The records distinguish deterministic mathematical guarantees, outward numerical enclosures and the independent-bin assumption supporting finite-sample inference.
\end{abstract}\end{frontmatter}
\section*{R61 preservation notice}
This is the retained development edition. The current R61 article and technical supplement contain the new exact-search and original-optimum verification results. The body below preserves its historical results, assumptions and observations.

\section{Primitive domain and exact moduli}\label{supp:primitive49}
Use the family in main equations~\eqref{M-eq:investment49}--\eqref{M-eq:cost49}, with $d=2,3,4$. Capacity is between $1/8$ and $1/4$ and has $\ell^1$ modulus $1/(8d)$. Every component of the dynamics is at least $1/32$. A direct upper bound is $29/32$ for the larger action exposure and $27/32$ for the smaller one. Thus the cube is invariant without an action-dependent exit event. The numerical clipping to $[0,1]$ is inactive on the true primitive image; it does not authorize out-of-domain evaluation.

For $d\ge2$, the sums of the absolute entries in a state column of the transition derivative are bounded by $1/2+1/16+1/8=11/16$. The action derivative has $\ell^1$ norm $\sum_j\gamma_j$. The action derivative of current cost is $2pa+16a^3\le p/2+1/4$. These are bounds between feasible pairs, not a claim that the same action is feasible at both states.

On either shortage region the state gradient is affine. The derivative of the squared-target and imbalance terms in coordinate $j$, after multiplication by $d$, is $4(x_j-5/8)+(2x_j-x_{j-1}-x_{j+1})/2$. In the shortage region add $-8(1/2-2\sum x/d)$. The terminal target derivative replaces 4 by 8. A linear functional on each region reaches its extrema at region vertices. For the declared dimensions, these are cube vertices and vertices with all but at most one coordinate at zero or one and $\sum x=d/4$. Exact evaluation gives absolute derivative bounds $13/(2d)$ and $9/d$. The added regression verifies these vertices rationally. The proof does not assert that these convenient formulas remain valid for untested arbitrary $d$.

The state covering radius of a nonuniform tensor grid is half the sum of its maximum coordinate spacings. Each action net is generated from a downward enclosure of capacity. Its coverage is at most $1/(8K)+2^{-48}$ under the implemented binary64 construction. This last conservative allowance covers capacity enclosure, multiplication, division and the final downward displacement on the bounded action domain. Fixed spacing and valid interval arithmetic are checked; a negative action is clipped to zero. Downward quantization on deployment preserves feasibility because the capacity itself is enclosed from below over the complete acquisition cell.

\section{Numerical labels and the separated account}\label{supp:numerical49}
An objective query has lower and upper endpoints for its midpoint-sum expectation, plus the deterministic continuous-law integration remainder. Labels are rounded on the inherited dyadic lattice, and the largest distance to the objective endpoints supplies the query allowance $e_t$. Selecting the least rounded label and retaining its original action costs one allowance in the label comparison and another in the complete selected-policy account. It does not require a second copy of the action-cover term in main equation~\eqref{M-eq:constructionloss49}.

For a uniform scalar bin of width $w$, the mean absolute distance to its midpoint is $w/4$. Here $w=1/(16M)$ and the transition's shock modulus is $d$. Therefore the numerical account decomposes into $\beta dL_{t+1}/(64M)$ and the remaining arithmetic width. The recorded primitive coefficients in main equation~\eqref{M-eq:coefficients49} reconstruct each witness bound as $a/N+b/K+q/M+\eta$ with a nonnegative numerical remainder. The release checks this decomposition at all witness rungs, as well as the full signed formula for every conventional rung.

The continuous allocation in main Proposition~\ref{M-prop:allocation49} minimizes a leading product, not a full finite-precision program. With positive coefficient $c_i$ and exponent $p_i$, changing variables to error contribution $x_i=c_i/r_i$ gives the objective $\sum_i p_i\log x_i$. Its constrained first-order condition is $x_i=\min(\bar x_i,p_i\lambda)$. The multiplier exists uniquely until every cap binds. A zero coefficient has no benefit from refinement, so its coordinate stays at the resource lower bound. Upward dyadic rounding keeps every error term weakly smaller. Compiler setup, integer counts, coefficient bit lengths, adaptive pilots, memory and durable output are not removed by that relaxation.

\section{Why the compiled actor is the original actor}\label{supp:owner49}
At a tensor node the distance transform minimizes over original sites, not merely neighboring transformed nodes. Each forward relaxation along a coordinate line compares the current value--owner pair with the preceding pair plus the weighted coordinate distance; the backward pass does the same from the opposite direction. The invariant is the lexicographic minimum over all processed original sites in the processed coordinates. Separability proves the invariant after the final coordinate. Keeping only the winning value would lose the actor identity.

For an off-grid point, choose a cell corner coordinatewise between the point and its lexicographically first minimizing site. Distances add exactly. At that corner an owner giving a smaller value, or a tied value and smaller original index, contradicts the minimizer at the original point. This proves the owner identity even for inconsistent labels, zero slope, boundary cells and ties. Nonuniform coordinate spacings enter only the exact distances.

In interval execution, value expressions are enclosed outward. Actor candidates whose lower objective exceeds another candidate's upper objective can be discarded; the remaining possibilities must be retained. An exact point tie is resolved by rational comparison and the original index rule. A genuine uncertain state box is not an exact point: the hull of its surviving actions is propagated. Boxes crossing cells use the union of all relevant corners. These operations may cost more than one point query, and the direct records retain actor and measurement ambiguity counts.

\section{Acquisition for witness and conventional policies}\label{supp:acquisition49}
Let $x\in C$ be the true state and $\widehat x$ the reported midpoint, with radius $r$ in $\ell^1$. Fix the selected original node $x_i$ and node action $a_i$. The ideal repair at the true state is feasible and within $K^A\|x-x_i\|_1$ of $a_i$. Replacing that repair by the cell-wide lower-capacity repair followed by downward spacing $\Delta$ changes the action by at most $2K^Ar+\Delta$. Thus
\[
 Q(x,a^{\rm impl})\le y_i+e+L^{\rm gr}\|x-x_i\|_1+ D(2K^Ar+\Delta).
\]
For a witness selection at $\widehat x$, $y_i+L^{\rm gr}\|\widehat x-x_i\|_1=f(\widehat x)$, up to a separately certified selector error. The triangle inequality and $L^f=L^{\rm gr}$ yield the allowance in main equation~\eqref{M-eq:acquire49}.

For conventional nearest-node selection, subtract and add $f(\widehat x)$. The expression $y_i+L^{\rm gr}\|\widehat x-x_i\|_1-f(\widehat x)$ is bounded by its complete nearest-cell maximum $E_t$. The remaining terms are at most $(L^{\rm gr}+L^f+2DK^A)r+D\Delta$. This proves main equation~\eqref{M-eq:fvisensor49} without a continuity assumption on the actor. It also explains why applying the witness coefficient unchanged to an interpolated conventional critic would be incorrect when the actual critic modulus differs.

For two-state records, $r(b)=2^{-b}$ and the common spacing is $1/4096$. The all-state bit criterion is therefore $G+K_\Delta+A2^{-b}+2\lambda b\sum_{t<T}\beta^t$, where $A$ uses the appropriate actual critic modulus and $K_\Delta=\sum_t\beta^tD_t/4096$. The interval-net criterion uses actual expected costs on the same finite grid instead of replacing $G+K_\Delta+A2^{-b}$ by a claim about true loss.

\section{Continuous path law and simultaneous inference}\label{supp:inference49}
Conditional on all construction records, the studied controllers are fixed. Each simulation row draws independent uniform indices for the initial coordinates and each innovation. The interval simulator encloses the complete continuous cell corresponding to those indices; no bin midpoint is substituted for the true path distribution in the estimand. Uniformly mixing conditional continuous cells reproduces the original uniform initial and innovation laws. Common innovations may couple two policies, but each follows its own endogenous trajectory.

In dimension two, a valid support interval for a raw discounted path-cost difference is $[-H,H]$, where
\[
 H=\frac{99+4p}{64}\sum_{t<T}\beta^t+\frac{37}{16}\beta^T.
\]
Indeed, current cost is nonnegative and at most $(99+4p)/64$ by separately bounding squared-target, imbalance, shortage and action costs. Terminal cost is at most $37/16$. These conservative bounds apply to every feasible repaired path. Clipping interval endpoints to the support preserves inclusion. Lower endpoint expectations lie below the true cost difference; upper endpoint expectations lie above it.

The stored moments consist of outward sample-mean intervals and upper sample-variance bounds for each endpoint sample. In \eqref{M-eq:radius49}, the square root is rounded upward after checking its square against the rational radicand. The linear remainder is rational. With $K\le300$, $\alpha=1/100$ and $\ell=13$, a finite positive Taylor sum verifies $e^{13}>4K/\alpha$. The bounded empirical-Bernstein inequalities of \citet{maurer2009} and a union bound therefore yield the claimed simultaneous event. Dependence between contrasts through shared economic objects is permitted; independence within each bin-index sample is the statistical assumption. The fixed pseudorandom seeds do not prove that assumption.

A first-crossing choice is measurable with respect to construction information. The conditional coverage assertion consequently remains valid after this choice. A bit selector or a fee queried after seeing the simultaneous intervals also needs no additional error allocation: on the one event, each underlying cost difference already belongs to its interval. This does not validate new initial laws, new trained policies, or unrecorded contrasts. Old and new study families have separate error accounts; combining them incurs the corresponding union bound.

\section{Exact frontiers and decision semantics}\label{supp:frontiers49}
For a fixed construction ladder, the first-crossing function can change only at a recorded certificate. A sorted union of all such values, together with zero, partitions all positive tolerances into left-closed, right-open intervals. At an attained endpoint the weak inequality in the stopping rule applies. Later nonmonotone or dominated certificates need not change the function, but remain in the raw record. Adjacent intervals with identical returned rungs may be merged for display; the exact rational partition is deposited.

For a direct interval $[l,u]$, the gain from replacing witness by FVI is in $[l,u]$; the reverse gain is in $[-u,-l]$. The fee $\max(0,u,-l)$ suffices to certify that neither switch is strictly profitable. The smaller $\max(0,l,-u)$ only says that neither switch has an interval lying wholly in its strictly profitable region. When zero lies inside a wide interval, this latter threshold is zero and is not evidence that a free switch cannot pay. The release reports both objects.

For information price $\lambda$, each bit choice supplies lower and upper affine net-cost lines. Enumerating all their nonnegative intersections partitions the price half-line. On an open interval, the active lower and upper lines are fixed. At a boundary all lines are compared exactly with smallest-bit ties. The difference between the active upper and lower lines is the selected implementation's valid regret upper bound. The identity line at 16 bits is exactly zero before information charges. No simulation count or time is assigned to an unperformed identity experiment; its arithmetic identity is recorded explicitly.

\section{Reference evaluation and safe improvement}\label{supp:improve50}
We give details of the reference-policy argument to separate a true improvement theorem from an approximate-greedy heuristic. All comparisons are on the same invariant domain and common conditional-evaluation domain as the main article. Feasible measurable policies include acquisition, capacity repair and numerical action realization.

Let $e_t=J_t^\pi-h_t$. At the terminal date $e_T\in[a_T,b_T]$. If $e_{t+1}\in[a_{t+1},b_{t+1}]$, cash invariance and monotonicity yield
\[
 \beta_ta_{t+1}\le
 \mathcal T_t^\pi J_{t+1}^\pi-\mathcal T_t^\pi h_{t+1}
 \le\beta_tb_{t+1}.
\]
Adding the signed own-policy residual proves the backward enclosure in Proposition~\ref{M-prop:advantage50}. Applying the same statement with either feasible action proves the difference allowance. No independence or differentiability of the fitted critic is required. The domain hypothesis matters: a certainty equivalent need not satisfy the argument for an inadmissible translated continuation.

If $\widetilde h_t=h_t+k_t$, then $\widetilde e_t=e_t-k_t$ and the residual shifts by $\beta_tk_{t+1}-k_t$. The propagated intervals become $[a_t-k_t,b_t-k_t]$, so their widths are unchanged. The two action evaluations of $\widetilde h_{t+1}$ each shift by $\beta_tk_{t+1}$, leaving the contrast unchanged. This is the centered invariance of the improvement test.

For the full expectation identity, define
$A_t^\pi(x,a)=Q_t(J_{t+1}^\pi;x,a)-J_t^\pi(x)$. Along any feasible $\mu$ trajectory,
\begin{align*}
 \E_x^\mu\sum_{t<T}B_tA_t^\pi(X_t,\mu_t(X_t))
 &=\E_x^\mu\sum_{t<T}\{B_tc_t(X_t,\mu_t(X_t))
          +B_{t+1}J_{t+1}^\pi(X_{t+1})-B_tJ_t^\pi(X_t)\}\\
 &=J_0^\mu(x)-J_0^\pi(x).
\end{align*}
Boundedness and finite horizon justify expectation and summation. Rejected actions give exactly zero reference advantage, while accepted ones give at most $-\kappa_t$. This proves Theorem~\ref{M-thm:improve50}. Under a nonlinear conditional recursion, use backward monotonicity instead of this telescoping expectation identity. A nonlinear evaluation does not justify the occupancy formula by changing notation alone.

The state-cell version fixes an action and verifies its inequality uniformly on the whole cell. Testing a cell midpoint alone is insufficient. Every implemented true state belongs to a cell whose gate is valid, and boundary cells use a declared deterministic report rule. In interval simulation, boxes straddling report boundaries retain all eligible actions, independently of whether their point-state selection is discontinuous. The analytic nonincrease may tighten a numerical advantage enclosure by intersection with $(-\infty,0]$, because the support is established independently of the simulated sample. This is not truncation of an unfavorable sample estimate.

\section{Exact integration and action optimization}\label{supp:terminal50}
Write $Z\sim U[-1/32,1/32]$. The deterministic part of each next coordinate is $y_j$, and its random part is $s_jZ$. Hence the expectation of $(y_j+s_jZ-c)^2$ is $(y_j-c)^2+1/3072$. For the cyclic difference use coefficient $s_j-s_{j+1}$. Summing yields the quadratic and variance terms in equation~\eqref{M-eq:exactterminal50}.

For the shortage term, the affine random component is uniform on $[-r,r]$, where $r=|\sum_js_j|/(16d)$. Direct antiderivatives give, for $r>0$,
\[
 \E(b+U)_+^2=\frac{(b+r)_+^3-(b-r)_+^3}{6r},\qquad
 \E(b+U)_+=\frac{(b+r)_+^2-(b-r)_+^2}{4r}.
\]
For $r=0$ take $b_+^2$ and $b_+$. These formulas hold at the two boundaries as well as inside the crossing region. The implementation encloses cancellation outward. Independent regression checks split the innovation interval at the exact shortage root and integrate the resulting quadratic pieces by exact rational Simpson identities; the checks are not used to generate the study's continuation labels.

At fixed $x$, the transition is affine in $a$. The terminal objective is a sum of convex quadratics and squared positive parts of affine functions. Its expectation is convex. The action cost $pa^2+4a^4$ has second derivative $2p+48a^2\ge2p$. Thus final conditional cost has a unique continuous minimizer, possibly on the capacity boundary. Its lattice minimizer lies at the first lattice point with nonnegative derivative or the preceding point, with boundary truncation. The implementation searches at most the fixed 4096 denominator's range, compares the two costs, and uses exact rational comparison when outward point intervals fail to decide a tie. The incumbent need not be one of these two actions: safety is established by the subsequent whole-cell comparison.

For bits $b$, let integer cell indices be $k_1,\ldots,k_d$. The smallest capacity on the cell is $1/8+\sum_jk_j/(8d2^b)$. Its largest lattice index is exactly
\[
 \left\lfloor\frac{4096(d2^b+\sum_jk_j)}{8d2^b}\right\rfloor.
\]
Integer arithmetic avoids an upward feasibility error, including in dimension three. The original action is clipped and rounded downward by this same rule. Feasibility is therefore not inferred from an interval midpoint. The final gate keeps the same earlier policies; its true improvement is relative to that acquired incumbent, not to a different exact-state policy.

\section{Simultaneous actual-cost inference}\label{supp:inference50}
The R50 primary family contains 25 pairs of baseline generators. Each group contains four policies---the two incumbents and their final repairs---and produces four absolute costs and six pairwise differences. Thus there are 250 sampled estimands. Their selection uses frozen construction information only. Within each group the continuous initial law and the first $T-1$ innovations are represented by forty-bit bins with outward propagation. The last innovation is integrated analytically. Conditional expectation over this last innovation preserves each policy's expected total cost. It reduces simulation variance without substituting a new innovation law.

The path-cost support derived in the retained inference section applies in dimensions two through four: the target and imbalance terms are coordinate averages with the same bounds, the shortage term is bounded by $1/2$, and feasible actions do not exceed $1/4$. Consequently $0\le J\text{-score}\le H$, with $H$ as in Section~\ref{supp:inference49}; a paired score lies in $[-H,H]$. An incumbent-minus-own-repair contrast is nonnegative by the independently proved gate and is bounded by $H$. Because the earlier trajectory is common, its numerical enclosure uses the centered final contrast instead of subtracting two complete total-cost intervals.

For endpoint observations $X_i\in[A,B]$, the retained empirical-Bernstein bound uses
\[
 R_n=\sqrt{2\overline{\operatorname{Var}}(X)\ell/n}
       +7(B-A)\ell/[3(n-1)],
\]
with outward moment intervals, an upper sample variance, and a rationally checked upward square root. Lower endpoints and upper endpoints each receive a one-sided bound. Set $n=131072$, $\ell=13$ and $\alpha=1/100$. The family cap is 512; an exact finite Taylor sum proves $e^{13}>4(512)/\alpha$. A union bound therefore covers all 250 actual-cost estimands under the independent-bin model. The cap leaves room but does not certify unrecorded models or initial laws. The pseudorandom stream and its hash establish reproducibility, not independence of ideal observations. The confidence event for the retained R49 family is separate; combining old and new statistical assertions incurs the sum of their failure probabilities.

\section{Mechanism and economic frontiers}\label{supp:mechanism50}
Denote witness and FVI by $W,F$ and their final repairs by $W^+,F^+$. Pure algebra yields
\begin{equation}\label{eq:decomposition50}
 J_W-J_F=(J_W-J_{W^+})+(J_{W^+}-J_{F^+})-(J_F-J_{F^+}).
\end{equation}
All three terms concern actual costs. The first and last are direct measures of removable final-action loss; the middle retains differences from earlier actions, acquired states, continuation approximation and the gate's conservatism. It is not identified as a single causal contribution of neural representation. The experiment reports each term, accepted and blocked cell counts, final action means, final state means and earlier action means by date. This exposes how much of the original ordering can be removed without retraining the continuation. It does not claim that the residual has been exhaustively decomposed into mutually independent mechanisms.

For finite intervals $[L_j,U_j]$, the upper-net-cost choice's true regret is bounded by the selected upper line minus the minimum lower line. To see why this is conservative, consider two overlapping cost intervals: their midpoint ordering may reverse within the same event. The upper-envelope decision remains valid but cannot identify the true minimizer. Pairwise centered intervals can certify individual replacements more sharply than the absolute frontier. For any fixed installation fee, transmission conversion factor and nonnegative resource price, add their declared costs before applying the comparison. A new price does not require new sampling on the same finite event. A new sensing policy does.

\section{Proofs for Finite-Sweep Accuracy}\label{app:finite-sweeps51}
\begin{proof}[Proof of Theorem~\ref{M-thm:finite-sweeps51}]
Fix a pass, date, and state, and suppress their indices. Write $J$ for the
incumbent continuation, $h$ for its approximator, $\pi$ for the incumbent
action, and $v$ for the proposal. By the signed evaluation band, for every
feasible action $a$,
\[
 \beta a_{t+1}\leq Q(J;x,a)-Q(h;x,a)\leq\beta b_{t+1}.
\]
Consequently, the incumbent-relative advantage of the proposal is at most
$\overline d+\beta w_{t+1}$. An accepted proposal therefore does not raise
the incumbent-continuation Bellman cost, while rejection retains exactly that
cost. Backward induction gives $J^{\pi^{k+1}}\leq J^{\pi^k}$ at every
date. This induction uses the new continuation for the implemented new
policy; the comparison itself uses the old continuation.

For an accepted proposal, approximate greediness and the band imply
\[
 Q(J;x,v)\leq (T_th)(x)+\alpha+\beta b_{t+1}
 \leq (T_tJ)(x)+\alpha+\beta w_{t+1}.
\]
For a rejected proposal, $\overline d+\beta w_{t+1}>0$ and
$\overline d\leq d+\zeta$. Thus
\[
 Q(h;x,\pi)<Q(h;x,v)+\zeta+\beta w_{t+1}
 \leq(T_th)(x)+\alpha+\zeta+\beta w_{t+1}.
\]
Transferring this inequality from $h$ to $J$ adds at most one further
$\beta w_{t+1}$. Both cases therefore yield
\begin{equation}\label{eq:approx-greedy-old51}
 T_t^{\pi^{k+1}}J_{t+1}^{\pi^k}
 \leq T_tJ_{t+1}^{\pi^k}+\epsilon_t^k.
\end{equation}
The already proved policy nonincrease, monotonicity of
$T_t^{\pi^{k+1}}$, and \eqref{eq:approx-greedy-old51} give
\[
 J_t^{\pi^{k+1}}\leq T_t^{\pi^{k+1}}J_{t+1}^{\pi^k}
 \leq T_tJ_{t+1}^{\pi^k}+\epsilon_t^k.
\]
The Bellman operator is $\beta_t$-Lipschitz in the supremum norm. Subtracting
$V_t^*=T_tV_{t+1}^*$ and using feasibility proves
\eqref{M-eq:sweep-recursion51}. Repeated substitution proves
\eqref{M-eq:sweep-unroll51}; at $m=T-t$ its terminal term vanishes.
Nonincrease of $E_t^k$ in further passes then gives
\eqref{M-eq:sweep-final51}. The exact case has zero error at every date.
\end{proof}

\begin{proof}[Proof of Proposition~\ref{M-prop:cell-greedy51}]
Fix $x\in C$ and any feasible $a\in A_t(x)$. Choose $m\in M_C$ with
$\|a-m\|\leq\delta_{A,C}$. Then
\begin{align*}
 Q_t(h;x,v_C)&\leq Q_t(h;x_C,v_C)+L_tr_C\\
 &\leq Q_t(h;x_C,m)+\eta_C+L_tr_C\\
 &\leq Q_t(h;x,a)+\eta_C+2L_tr_C+D_t\delta_{A,C}.
\end{align*}
Taking an infimum over $a$ proves the assertion without assuming that the
infimum is attained. Menu feasibility is uniform on $C$ by hypothesis.
\end{proof}

\begin{proof}[Proof of Proposition~\ref{M-prop:sharp-width51}]
Take $T=1$, zero stage cost, $\beta_0=1$, and two deterministic actions
leading to different terminal states $y_0,y_1$. Let $w>0$ and
$0<\delta<w$. Set
\[
 h_1(y_0)=w,\quad h_1(y_1)=\delta,\qquad
 g(y_0)=2w,\quad g(y_1)=\delta.
\]
The exact terminal error lies in $[0,w]$. The incumbent chooses $y_0$ and
the exact $h_1$-greedy proposal chooses $y_1$. Its exact contrast is
$d=\delta-w$, with $\alpha=\zeta=0$. The safe upper comparison is
$d+w=\delta>0$, so the gate retains the incumbent. Its actual optimality
gap is $2w-\delta$. Sending $\delta$ down to zero establishes sharpness.
This example also shows why rejecting an uncertified proposal is not evidence
that the proposal lacks a true economic benefit.
\end{proof}

\paragraph{Computed verification and logical scope.}
The source-bound exact-rational tests use finite transition matrices, signed
residual propagation, and independently evaluated policy and optimal costs.
They verify policy nonincrease, both gate cases, the one-pass recurrence,
constant-shift invariance, finite-horizon exact termination, and the two-action
sharpness construction. They are regression tests for the implementation of
the displayed formulas; the preceding arguments, rather than enumeration,
prove the theorems for the stated model class. When bands are statistical,
all passes and proposed comparisons require a common simultaneous coverage
event or a valid sequential error allocation. Reusing a confidence interval
outside its stated family does not satisfy the hypotheses.

\section{Proofs for policy-relevant continuation errors}\label{supp:contrast52}
\begin{proof}[Proof of Theorem~\ref{M-thm:contrast52}]
Fix a pass and suppress its index. For any feasible $a,b$, subtraction of the two Bellman expressions gives
\[
 Q_t(J_{t+1}^\pi;x,a)-Q_t(J_{t+1}^\pi;x,b)
 =Q_t(h_{t+1};x,a)-Q_t(h_{t+1};x,b)
       +\beta_t(P_t^a-P_t^b)e_{t+1}(x).
\]
Thus the accepted proposal has nonpositive true advantage relative to the incumbent. At a rejected state the implemented action is exactly the incumbent, whose advantage is zero. Backward induction gives $J^{\pi^+}\le J^\pi$. This proves safety even when only the proposal--incumbent contrast is certified.

For accuracy, suppose first that the proposal is accepted. For every feasible $a$, its approximate greediness and the uniform pair certificate imply
\[
 Q_t(J_{t+1}^\pi;x,v)-Q_t(J_{t+1}^\pi;x,a)
 \le\alpha_t+\beta_t\chi_t.
\]
Taking the infimum over $a$ requires neither existence of a minimizer nor a new selection assumption. If the proposal is rejected, then
$\overline d_t+\beta_tR_t(x;v,\pi_t)>0$ and
$\overline d_t\le d_t+\zeta_t$. Hence
\[
 Q_t(h_{t+1};x,\pi_t)
 <Q_t(h_{t+1};x,v)+\zeta_t+\beta_tR_t(x;v,\pi_t).
\]
Comparison with each feasible $a$ and another use of the action certificate yield an excess over $T_tJ_{t+1}^\pi$ of at most
$\alpha_t+\zeta_t+2\beta_t\chi_t$.
This common upper bound covers acceptance as well. Since future costs weakly decrease,
\[
 J_t^{\pi^+}=\mathcal T_t^{\pi^+}J_{t+1}^{\pi^+}
 \le\mathcal T_t^{\pi^+}J_{t+1}^{\pi}
 \le T_tJ_{t+1}^{\pi}+\tilde\epsilon_t.
\]
Monotonicity and the $\beta_t$ supremum-norm Lipschitz property of $T_t$ give the asserted recursion. The terminal gap is zero, and repeated substitution along successive passes gives the stated finite-horizon unrolling. Every retained policy is feasible for the original optimum, so a zero upper bound implies equality with that optimum. A fixed information contract may preclude zero greedy-search error; the theorem does not assume it away.

For sharpness take one period, two actions leading to distinct deterministic terminal states, zero current cost, and $\beta=1$. For $w>0$ and $0<\delta<w$, let the terminal critic be $(w,\delta)$ and the terminal cost be $(2w,\delta)$. The incumbent selects the first state and the exact greedy proposal the second. The exact action-error contrast is $w$, so $\chi=w$. The gate rejects because the critic contrast plus its allowance is $\delta>0$. The retained policy's gap is $2w-\delta$, whereas the terminal optimality gap and both search and arithmetic errors are zero. Letting $\delta$ decrease to zero excludes every universal coefficient below two.
\end{proof}

\begin{proof}[Proof of Proposition~\ref{M-prop:coupled52}]
The terminal assertion follows from $e_T=r_T$. At an earlier date, the policy-evaluation identity is
$e_t(x)=r_t(x)+\beta_tP_t^{\pi_t(x)}e_{t+1}(x)$.
The coupling's two marginals allow subtraction at states $x,y$ inside one integral. The triangle inequality and the inductive pair bound give
\[
 |e_t(x)-e_t(y)|\le H_t(x,y)+\beta_t
       \int D_{t+1}(u,v)\,d\Lambda_t^\pi(x,y)(u,v).
\]
The signed range supplies the separate bound $w_t$, so the minimum is valid. Coupling two action kernels at the same state gives the identical argument without the residual term and establishes the action certificate. Replacing an integral by an upper enclosure can only enlarge the bound before its intersection with the separately valid scalar width. This argument uses the marginals, not independence or optimality of the coupling.
\end{proof}

\begin{proof}[Proof of Corollary~\ref{M-cor:null52}]
The action difference annihilates $n$. Each of the two conditional means of $u$ lies in its verified interval, hence their difference in absolute value is at most its width. Further,
$J^\pi-(h+n)=u$, so the existing signed-residual recursion applied to the corrected critic yields the required interval without using an unknown true value. Adding $n$ to the continuation changes every feasible action objective by the same state-dependent amount; all action contrasts and approximate-greedy inequalities are therefore unchanged. The conclusion applies separately at each date and pass for which the membership and residual certificates are supplied.
\end{proof}

\begin{proof}[Proof of Proposition~\ref{M-prop:investment-null52}]
In the original dynamics the coefficients of $a$ in the first two coordinates are $1/2$ and $1/4$. For each common innovation $z$, the coefficient in $F_1(x,a,z)-2F_2(x,a,z)$ is consequently zero. Every function of that difference has identical realized values for two feasible actions with the same $z$, and thus identical conditional means under the original uniform innovation law. On the cube $y_1-2y_2+1$ ranges from $-1$ to $2$; its positive part ranges from zero to two. The global terminal error width is exactly $2M$. Its action contrast is nevertheless zero. The corrected critic is the original terminal cost $g$, for which the evaluation error vanishes. Proposal contrasts and their symbolic, outward-enclosed implementation are unchanged. Keeping the same proposals, tie conventions and cellwise gate therefore gives exactly the same terminal policy for every amplitude, not merely a similar observed cost.
\end{proof}

\subsection{Finite-state reference and limitations}
The exact reference implementation constructs a common-uniform coupling by exhausting ordered probability masses. It verifies both marginals, propagates pairwise residual bounds backward, and forms every same-state action contrast. Its regression suite independently evaluates the true finite-state policy values to test the certificates; those true values are not used by the certificate constructor. Tests cover one through five full improvement passes, a nonconstant action-null error with amplitudes zero, one, sixteen and 4096, non-null errors, constant shifts, exact feasibility, acceptance and rejection, and the sharp two-action example. These finite-state programs validate executable identities and counterexamples; they do not discretize the investment benchmark or establish full-horizon empirical convergence in that continuous economy.

\section{Complete construction records}\label{supp:records49}
The next table reports all distinct method, economy and resource rungs. Each row has three isolated repetitions with identical policy and certificate checkpoint hashes. Times are the median complete prefix, followed by its minimum and maximum; they do not omit unsuccessful earlier rungs. Counts and storage at full precision are in the machine-readable records. The table's decimal rounding is descriptive: exact rational bounds, not printed rounding, determine attainment.
{\small
\setlength{\tabcolsep}{3.5pt}
\begin{longtable}{@{}llrrrr@{}}
\caption{All distinct construction rungs and complete-prefix times}\label{tab:rungs49}\\
\toprule
$d;T;p$ & Method & $N,K,M$ & Bound & Median sec. & Range sec. \\
\midrule
\endfirsthead
\toprule
$d;T;p$ & Method & $N,K,M$ & Bound & Median sec. & Range sec. \\
\midrule
\endhead
\bottomrule
\endfoot
2;2;1 & W & 4,2,1 & 10.0042 & 0.014 & [0.014,0.017] \\
2;2;1 & W & 8,2,1 & 5.591673 & 0.034 & [0.034,0.036] \\
2;2;1 & W & 16,4,2 & 2.795836 & 0.092 & [0.091,0.093] \\
2;2;1 & W & 32,4,2 & 1.692705 & 0.290 & [0.290,0.291] \\
2;2;1 & W & 64,8,4 & 0.8463524 & 1.458 & [1.449,1.459] \\
2;2;1 & W & 128,16,8 & 0.4231762 & 10.658 & [10.612,10.696] \\
2;2;1 & A & 4,2,1 & 10.68755 & 0.042 & [0.042,0.043] \\
2;2;1 & A & 8,2,1 & 6.132349 & 0.090 & [0.090,0.091] \\
2;2;1 & A & 16,4,2 & 3.216146 & 0.176 & [0.176,0.178] \\
2;2;1 & A & 32,4,2 & 1.906506 & 0.399 & [0.398,0.402] \\
2;2;1 & A & 64,8,4 & 0.9622566 & 1.514 & [1.509,1.521] \\
2;2;1 & A & 128,16,8 & 0.4833928 & 9.817 & [9.748,9.900] \\
2;2;1 & F & 4,2,1 & 10.68755 & 0.029 & [0.029,0.029] \\
2;2;1 & F & 8,2,1 & 6.132349 & 0.061 & [0.061,0.063] \\
2;2;1 & F & 16,4,2 & 3.216146 & 0.119 & [0.118,0.122] \\
2;2;1 & F & 32,4,2 & 1.906506 & 0.269 & [0.268,0.274] \\
2;2;1 & F & 64,8,4 & 0.9622566 & 1.134 & [1.130,1.143] \\
2;2;1 & F & 128,16,8 & 0.4833928 & 8.458 & [8.455,8.491] \\
2;2;4 & W & 4,2,1 & 10.31855 & 0.014 & [0.014,0.014] \\
2;2;4 & W & 8,2,1 & 5.844473 & 0.034 & [0.034,0.035] \\
2;2;4 & W & 16,4,2 & 2.922236 & 0.092 & [0.092,0.096] \\
2;2;4 & W & 32,4,2 & 1.803718 & 0.290 & [0.287,0.303] \\
2;2;4 & W & 64,8,4 & 0.9018591 & 1.452 & [1.452,1.510] \\
2;2;4 & W & 128,16,8 & 0.4509295 & 10.689 & [10.584,10.884] \\
2;2;4 & A & 4,2,1 & 10.92323 & 0.043 & [0.042,0.043] \\
2;2;4 & A & 8,2,1 & 6.330362 & 0.091 & [0.090,0.092] \\
2;2;4 & A & 16,4,2 & 3.31506 & 0.177 & [0.176,0.179] \\
2;2;4 & A & 32,4,2 & 2.000228 & 0.397 & [0.397,0.403] \\
2;2;4 & A & 64,8,4 & 1.009136 & 1.509 & [1.509,1.524] \\
2;2;4 & A & 128,16,8 & 0.5068371 & 9.797 & [9.763,9.822] \\
2;2;4 & F & 4,2,1 & 10.92323 & 0.029 & [0.028,0.033] \\
2;2;4 & F & 8,2,1 & 6.330362 & 0.061 & [0.060,0.065] \\
2;2;4 & F & 16,4,2 & 3.31506 & 0.118 & [0.118,0.123] \\
2;2;4 & F & 32,4,2 & 2.000228 & 0.268 & [0.268,0.273] \\
2;2;4 & F & 64,8,4 & 1.009136 & 1.134 & [1.131,1.145] \\
2;2;4 & F & 128,16,8 & 0.5068371 & 8.453 & [8.373,8.520] \\
2;3;1 & W & 4,2,1 & 14.5143 & 0.019 & [0.019,0.019] \\
2;3;1 & W & 8,2,1 & 8.228071 & 0.046 & [0.046,0.047] \\
2;3;1 & W & 16,4,2 & 4.114036 & 0.126 & [0.126,0.129] \\
2;3;1 & W & 32,4,2 & 2.542479 & 0.401 & [0.398,0.408] \\
2;3;1 & W & 64,8,4 & 1.271239 & 2.070 & [2.062,2.110] \\
2;3;1 & W & 128,16,8 & 0.6356197 & 15.635 & [15.607,15.924] \\
2;3;1 & A & 4,2,1 & 14.97615 & 0.061 & [0.060,0.062] \\
2;3;1 & A & 8,2,1 & 8.615403 & 0.130 & [0.129,0.131] \\
2;3;1 & A & 16,4,2 & 4.533221 & 0.252 & [0.251,0.255] \\
2;3;1 & A & 32,4,2 & 2.723572 & 0.568 & [0.565,0.574] \\
2;3;1 & A & 64,8,4 & 1.374878 & 2.148 & [2.145,2.184] \\
2;3;1 & A & 128,16,8 & 0.6908084 & 14.176 & [14.130,14.444] \\
2;3;1 & F & 4,2,1 & 14.97615 & 0.041 & [0.041,0.042] \\
2;3;1 & F & 8,2,1 & 8.615403 & 0.088 & [0.087,0.088] \\
2;3;1 & F & 16,4,2 & 4.533221 & 0.171 & [0.170,0.192] \\
2;3;1 & F & 32,4,2 & 2.723572 & 0.388 & [0.388,0.412] \\
2;3;1 & F & 64,8,4 & 1.374878 & 1.651 & [1.647,1.677] \\
2;3;1 & F & 128,16,8 & 0.6908084 & 12.420 & [12.409,12.456] \\
2;3;4 & W & 4,2,1 & 15.02035 & 0.019 & [0.019,0.019] \\
2;3;4 & W & 8,2,1 & 8.625736 & 0.046 & [0.046,0.048] \\
2;3;4 & W & 16,4,2 & 4.312868 & 0.127 & [0.125,0.129] \\
2;3;4 & W & 32,4,2 & 2.714215 & 0.402 & [0.397,0.408] \\
2;3;4 & W & 64,8,4 & 1.357108 & 2.085 & [2.053,2.103] \\
2;3;4 & W & 128,16,8 & 0.6785539 & 15.815 & [15.506,15.955] \\
2;3;4 & A & 4,2,1 & 15.30603 & 0.061 & [0.061,0.061] \\
2;3;4 & A & 8,2,1 & 8.884068 & 0.131 & [0.130,0.131] \\
2;3;4 & A & 16,4,2 & 4.667275 & 0.253 & [0.252,0.254] \\
2;3;4 & A & 32,4,2 & 2.853489 & 0.569 & [0.569,0.569] \\
2;3;4 & A & 64,8,4 & 1.440006 & 2.169 & [2.165,2.170] \\
2;3;4 & A & 128,16,8 & 0.723386 & 14.257 & [14.181,14.260] \\
2;3;4 & F & 4,2,1 & 15.30603 & 0.042 & [0.041,0.042] \\
2;3;4 & F & 8,2,1 & 8.884068 & 0.088 & [0.087,0.088] \\
2;3;4 & F & 16,4,2 & 4.667275 & 0.170 & [0.170,0.171] \\
2;3;4 & F & 32,4,2 & 2.853489 & 0.390 & [0.387,0.391] \\
2;3;4 & F & 64,8,4 & 1.440006 & 1.643 & [1.637,1.647] \\
2;3;4 & F & 128,16,8 & 0.723386 & 12.435 & [12.393,12.446] \\
3;2;1 & W & 4,2,1 & 10.09191 & 0.033 & [0.033,0.033] \\
3;2;1 & W & 8,4,2 & 5.045957 & 0.225 & [0.224,0.225] \\
3;2;1 & W & 16,8,4 & 2.522979 & 2.424 & [2.422,2.453] \\
3;2;1 & W & 32,8,4 & 1.415656 & 16.615 & [16.590,16.644] \\
3;2;1 & F & 4,2,1 & 10.91291 & 0.096 & [0.095,0.096] \\
3;2;1 & F & 8,4,2 & 5.756259 & 0.318 & [0.318,0.322] \\
3;2;1 & F & 16,8,4 & 2.983083 & 2.074 & [2.059,2.097] \\
3;2;1 & F & 32,8,4 & 1.654898 & 13.687 & [13.539,13.776] \\
3;3;1 & W & 4,2,1 & 14.67402 & 0.045 & [0.044,0.046] \\
3;3;1 & W & 8,4,2 & 7.337012 & 0.317 & [0.313,0.322] \\
3;3;1 & W & 16,8,4 & 3.668506 & 3.508 & [3.476,3.527] \\
3;3;1 & W & 32,8,4 & 2.088657 & 23.987 & [23.952,24.055] \\
3;3;1 & F & 4,2,1 & 15.28673 & 0.140 & [0.140,0.141] \\
3;3;1 & F & 8,4,2 & 8.038961 & 0.467 & [0.463,0.468] \\
3;3;1 & F & 16,8,4 & 4.173215 & 3.032 & [3.027,3.059] \\
3;3;1 & F & 32,8,4 & 2.334703 & 20.040 & [19.984,20.139] \\
4;2;1 & W & 4,2,1 & 10.0042 & 0.167 & [0.164,0.168] \\
4;2;1 & W & 8,4,2 & 5.002099 & 2.588 & [2.561,2.635] \\
4;2;1 & W & 16,8,4 & 2.50105 & 58.569 & [58.203,59.507] \\
4;2;1 & F & 4,2,1 & 10.98489 & 0.523 & [0.520,0.526] \\
4;2;1 & F & 8,4,2 & 5.772966 & 4.388 & [4.386,4.449] \\
4;2;1 & F & 16,8,4 & 2.981482 & 60.809 & [60.637,61.583] \\
4;3;1 & W & 4,2,1 & 14.5143 & 0.230 & [0.230,0.231] \\
4;3;1 & W & 8,4,2 & 7.257149 & 3.663 & [3.649,3.688] \\
4;3;1 & W & 16,8,4 & 3.628575 & 82.156 & [81.795,82.676] \\
4;3;1 & F & 4,2,1 & 15.40485 & 0.782 & [0.776,0.784] \\
4;3;1 & F & 8,4,2 & 8.075833 & 6.457 & [6.444,6.547] \\
4;3;1 & F & 16,8,4 & 4.176837 & 90.099 & [89.833,90.155] \\
\end{longtable}
}


\section{All direct returned-policy comparisons}\label{supp:direct49}
The labels U, L, C and H denote the uniform initial law and point laws at $(1/8,1/8)$, $(1/2,1/2)$ and $(7/8,7/8)$. Sensor E denotes full information; sensor 8 uses the eight-bit cell contract and downward action spacing. The historical favorable pair is marked R48; other pairs are the actual R49 first crossings. Every interval concerns expected witness cost minus expected FVI cost. All sampled intervals in this table and the next belong to the same new 99 percent family under the independent-bin model.
{\small
\setlength{\tabcolsep}{3.5pt}
\begin{longtable}{@{}llrrrrr@{}}
\caption{Every actual first-crossing policy-cost interval}\label{tab:all-direct49}\\
\toprule
Pair & $T;p$ & $\varepsilon$ & Law & Bits & Lower & Upper \\
\midrule
\endfirsthead
\toprule
Pair & $T;p$ & $\varepsilon$ & Law & Bits & Lower & Upper \\
\midrule
\endhead
\bottomrule
\endfoot
R49 & 2;1 & 4 & U & E & 0.00543396 & 0.00798729 \\
R49 & 2;1 & 4 & U & 8 & 0.00541635 & 0.00796936 \\
R49 & 2;1 & 4 & L & E & -0.00156041 & 0.000833084 \\
R49 & 2;1 & 4 & L & 8 & -0.00156513 & 0.000828396 \\
R49 & 2;1 & 4 & C & E & 0.00372573 & 0.0061508 \\
R49 & 2;1 & 4 & C & 8 & 0.00375513 & 0.00618077 \\
R49 & 2;1 & 4 & H & E & 0.00609378 & 0.00853627 \\
R49 & 2;1 & 4 & H & 8 & 0.00589939 & 0.00833643 \\
R49 & 2;1 & 2 & U & E & 0.0143531 & 0.0170811 \\
R49 & 2;1 & 2 & U & 8 & 0.0143519 & 0.0170783 \\
R49 & 2;1 & 2 & L & E & 0.0547607 & 0.0574082 \\
R49 & 2;1 & 2 & L & 8 & 0.0562672 & 0.0589226 \\
R49 & 2;1 & 2 & C & E & 0.00754534 & 0.0100938 \\
R49 & 2;1 & 2 & C & 8 & 0.00764868 & 0.0101989 \\
R49 & 2;1 & 2 & H & E & 0.00545776 & 0.00790326 \\
R49 & 2;1 & 2 & H & 8 & 0.00540822 & 0.00785334 \\
R49 & 2;1 & 1 & U & E & 0.00608323 & 0.00862331 \\
R49 & 2;1 & 1 & U & 8 & 0.00608821 & 0.00862733 \\
R49 & 2;1 & 1 & L & E & 0.0241465 & 0.0266462 \\
R49 & 2;1 & 1 & L & 8 & 0.0243754 & 0.0268759 \\
R49 & 2;1 & 1 & C & E & 0.0101329 & 0.0126139 \\
R49 & 2;1 & 1 & C & 8 & 0.0101188 & 0.012601 \\
R49 & 2;1 & 1 & H & E & 0.00787923 & 0.0103208 \\
R49 & 2;1 & 1 & H & 8 & 0.00773077 & 0.0101721 \\
R49 & 2;4 & 4 & U & E & 0.0048609 & 0.00758426 \\
R49 & 2;4 & 4 & U & 8 & 0.00488276 & 0.00760722 \\
R49 & 2;4 & 4 & L & E & 0.0185747 & 0.0211388 \\
R49 & 2;4 & 4 & L & 8 & 0.0186476 & 0.0212113 \\
R49 & 2;4 & 4 & C & E & 0.0347841 & 0.0374079 \\
R49 & 2;4 & 4 & C & 8 & 0.0346533 & 0.0372768 \\
R49 & 2;4 & 4 & H & E & -8.33958e-05 & 0.00253601 \\
R49 & 2;4 & 4 & H & 8 & 0.000139912 & 0.00276397 \\
R49 & 2;4 & 2 & U & E & 0.0126061 & 0.0154244 \\
R49 & 2;4 & 2 & U & 8 & 0.0124777 & 0.0152934 \\
R49 & 2;4 & 2 & L & E & 0.0213281 & 0.0239062 \\
R49 & 2;4 & 2 & L & 8 & 0.0216307 & 0.0242254 \\
R49 & 2;4 & 2 & C & E & 0.0363241 & 0.0389543 \\
R49 & 2;4 & 2 & C & 8 & 0.0365806 & 0.0392119 \\
R49 & 2;4 & 2 & H & E & 0.0408571 & 0.0435165 \\
R49 & 2;4 & 2 & H & 8 & 0.0405602 & 0.0432247 \\
R49 & 2;4 & 1 & U & E & 0.00493422 & 0.0076087 \\
R49 & 2;4 & 1 & U & 8 & 0.00498356 & 0.00765916 \\
R49 & 2;4 & 1 & L & E & 0.00791584 & 0.0104887 \\
R49 & 2;4 & 1 & L & 8 & 0.00799778 & 0.010571 \\
R49 & 2;4 & 1 & C & E & 0.016087 & 0.018698 \\
R49 & 2;4 & 1 & C & 8 & 0.0157071 & 0.0183196 \\
R49 & 2;4 & 1 & H & E & 0.00782793 & 0.0104847 \\
R49 & 2;4 & 1 & H & 8 & 0.00787666 & 0.0105306 \\
R49 & 3;1 & 4 & U & E & 0.0253656 & 0.0288147 \\
R49 & 3;1 & 4 & U & 8 & 0.0252662 & 0.0287151 \\
R49 & 3;1 & 4 & L & E & 0.0788767 & 0.0824397 \\
R49 & 3;1 & 4 & L & 8 & 0.0778317 & 0.0813824 \\
R49 & 3;1 & 4 & C & E & 0.100968 & 0.104201 \\
R49 & 3;1 & 4 & C & 8 & 0.101486 & 0.104716 \\
R49 & 3;1 & 4 & H & E & 0.0197496 & 0.0229893 \\
R49 & 3;1 & 4 & H & 8 & 0.019754 & 0.0229886 \\
R49 & 3;1 & 2 & U & E & 0.0109993 & 0.0141923 \\
R49 & 3;1 & 2 & U & 8 & 0.0110585 & 0.0142512 \\
R49 & 3;1 & 2 & L & E & 0.0282946 & 0.0314077 \\
R49 & 3;1 & 2 & L & 8 & 0.0279109 & 0.0310247 \\
R49 & 3;1 & 2 & C & E & 0.0176355 & 0.0207641 \\
R49 & 3;1 & 2 & C & 8 & 0.0171906 & 0.0203209 \\
R49 & 3;1 & 2 & H & E & 0.0126606 & 0.0156975 \\
R49 & 3;1 & 2 & H & 8 & 0.0126999 & 0.0157435 \\
R49 & 3;1 & 1 & U & E & 0.00423544 & 0.00731756 \\
R49 & 3;1 & 1 & U & 8 & 0.00451493 & 0.0075994 \\
R49 & 3;1 & 1 & L & E & 0.00514129 & 0.008227 \\
R49 & 3;1 & 1 & L & 8 & 0.00545945 & 0.00854446 \\
R49 & 3;1 & 1 & C & E & 0.00460682 & 0.00763251 \\
R49 & 3;1 & 1 & C & 8 & 0.00492024 & 0.00794768 \\
R49 & 3;1 & 1 & H & E & 0.00035797 & 0.00337421 \\
R49 & 3;1 & 1 & H & 8 & 0.00036318 & 0.00337929 \\
R49 & 3;4 & 4 & U & E & 0.0186753 & 0.022252 \\
R49 & 3;4 & 4 & U & 8 & 0.0186419 & 0.0222187 \\
R49 & 3;4 & 4 & L & E & 0.0276014 & 0.0309773 \\
R49 & 3;4 & 4 & L & 8 & 0.0279277 & 0.0313061 \\
R49 & 3;4 & 4 & C & E & 0.0549145 & 0.058394 \\
R49 & 3;4 & 4 & C & 8 & 0.0535423 & 0.0569789 \\
R49 & 3;4 & 4 & H & E & 0.0500905 & 0.0536733 \\
R49 & 3;4 & 4 & H & 8 & 0.0501068 & 0.0536738 \\
R49 & 3;4 & 2 & U & E & 0.00722022 & 0.0105987 \\
R49 & 3;4 & 2 & U & 8 & 0.00727216 & 0.0106525 \\
R49 & 3;4 & 2 & L & E & 0.0111161 & 0.0144007 \\
R49 & 3;4 & 2 & L & 8 & 0.0111398 & 0.0144261 \\
R49 & 3;4 & 2 & C & E & 0.0119992 & 0.0152703 \\
R49 & 3;4 & 2 & C & 8 & 0.0120479 & 0.0153198 \\
R49 & 3;4 & 2 & H & E & 0.00896933 & 0.0122714 \\
R49 & 3;4 & 2 & H & 8 & 0.00896107 & 0.0122625 \\
R49 & 3;4 & 1 & U & E & 0.00231571 & 0.00560954 \\
R49 & 3;4 & 1 & U & 8 & 0.00249689 & 0.00579256 \\
R49 & 3;4 & 1 & L & E & 0.00405598 & 0.00731121 \\
R49 & 3;4 & 1 & L & 8 & 0.000239139 & 0.00350007 \\
R49 & 3;4 & 1 & C & E & 0.00285917 & 0.00611809 \\
R49 & 3;4 & 1 & C & 8 & 0.00270615 & 0.00596212 \\
R49 & 3;4 & 1 & H & E & 0.000774402 & 0.00403881 \\
R49 & 3;4 & 1 & H & 8 & 0.000787911 & 0.0040506 \\
R48 & 3;4 & 2 & U & E & -0.00139342 & 0.00183546 \\
R48 & 3;4 & 2 & U & 8 & -0.00138819 & 0.00184069 \\
R48 & 3;4 & 2 & L & E & -0.00148104 & 0.00174625 \\
R48 & 3;4 & 2 & L & 8 & -0.00148102 & 0.00174626 \\
R48 & 3;4 & 2 & C & E & -0.00135341 & 0.0018749 \\
R48 & 3;4 & 2 & C & 8 & -0.00134267 & 0.00188569 \\
R48 & 3;4 & 2 & H & E & -0.00139997 & 0.00182783 \\
R48 & 3;4 & 2 & H & 8 & -0.00141355 & 0.00181414 \\
\end{longtable}
}
\noindent{\footnotesize Endpoints are displayed to six significant digits; exact rational endpoints and policy hashes govern every sign and fee conclusion. R48 rows evaluate the old favorable N=32 witness versus N=64 FVI pair under fresh independent interval paths.\par}


\section{All sensor-policy contrasts}\label{supp:sensors49}
Each row compares the actual $b$-bit implementation with the same returned controller at 16 bits, under the uniform initial law. The 16-bit self-rows are identities. The corresponding price selectors and exact continuous-price envelopes use these rows without additional sampling. Upper-net-cost regret is relative to the best bit in this catalogue, not an unobserved optimum over continuous precision or arbitrary sensors.
{\small
\setlength{\tabcolsep}{3.5pt}
\begin{longtable}{@{}llrrr@{}}
\caption{Every actual sensor-cost difference relative to 16 bits}\label{tab:all-sensor49}\\
\toprule
$T;p$ & Method & Bits & Lower & Upper \\
\midrule
\endfirsthead
\toprule
$T;p$ & Method & Bits & Lower & Upper \\
\midrule
\endhead
\bottomrule
\endfoot
2;1 & W & 4 & 0.000261375 & 0.00285711 \\
2;1 & W & 5 & -0.000427494 & 0.00215671 \\
2;1 & W & 6 & -0.00204196 & 0.000521085 \\
2;1 & W & 7 & -0.00088098 & 0.00160987 \\
2;1 & W & 8 & -0.00122806 & 0.00124059 \\
2;1 & W & 9 & -0.00121323 & 0.00123098 \\
2;1 & W & 10 & -0.00120642 & 0.00122091 \\
2;1 & W & 11 & -0.00119893 & 0.00121541 \\
2;1 & W & 12 & -0.00120501 & 0.00120022 \\
2;1 & W & 13 & -0.0011989 & 0.00119995 \\
2;1 & W & 14 & -0.00119712 & 0.00119717 \\
2;1 & W & 15 & -0.00119517 & 0.00119587 \\
2;1 & W & 16 & 0 & 0 \\
2;1 & F & 4 & 0.000199127 & 0.0026252 \\
2;1 & F & 5 & -0.000584413 & 0.00181963 \\
2;1 & F & 6 & -0.00085593 & 0.00153971 \\
2;1 & F & 7 & -0.00111845 & 0.00126841 \\
2;1 & F & 8 & -0.00117522 & 0.00120999 \\
2;1 & F & 9 & -0.00118398 & 0.00120093 \\
2;1 & F & 10 & -0.00118806 & 0.00119669 \\
2;1 & F & 11 & -0.00119023 & 0.00119442 \\
2;1 & F & 12 & -0.00119128 & 0.00119331 \\
2;1 & F & 13 & -0.00119182 & 0.00119274 \\
2;1 & F & 14 & -0.00119207 & 0.00119247 \\
2;1 & F & 15 & -0.00119219 & 0.00119235 \\
2;1 & F & 16 & 0 & 0 \\
2;4 & W & 4 & -4.77011e-05 & 0.00267196 \\
2;4 & W & 5 & -0.000905067 & 0.00180125 \\
2;4 & W & 6 & -0.00220586 & 0.000485462 \\
2;4 & W & 7 & -0.00108854 & 0.00153925 \\
2;4 & W & 8 & -0.00124312 & 0.00137376 \\
2;4 & W & 9 & -0.00131086 & 0.00128822 \\
2;4 & W & 10 & -0.00128939 & 0.00129611 \\
2;4 & W & 11 & -0.00128416 & 0.00129219 \\
2;4 & W & 12 & -0.00128539 & 0.00128378 \\
2;4 & W & 13 & -0.00128173 & 0.0012827 \\
2;4 & W & 14 & -0.00128018 & 0.00128022 \\
2;4 & W & 15 & -0.00127809 & 0.00128043 \\
2;4 & W & 16 & 0 & 0 \\
2;4 & F & 4 & -0.00112704 & 0.00143668 \\
2;4 & F & 5 & -0.00121844 & 0.00133951 \\
2;4 & F & 6 & -0.00125571 & 0.00129926 \\
2;4 & F & 7 & -0.00126208 & 0.00129244 \\
2;4 & F & 8 & -0.00127554 & 0.00127722 \\
2;4 & F & 9 & -0.00127593 & 0.00127681 \\
2;4 & F & 10 & -0.00127615 & 0.00127658 \\
2;4 & F & 11 & -0.00127626 & 0.00127646 \\
2;4 & F & 12 & -0.0012763 & 0.00127642 \\
2;4 & F & 13 & -0.00127634 & 0.00127638 \\
2;4 & F & 14 & -0.00127635 & 0.00127637 \\
2;4 & F & 15 & -0.00127636 & 0.00127636 \\
2;4 & F & 16 & 0 & 0 \\
3;1 & W & 4 & -0.000101036 & 0.00300605 \\
3;1 & W & 5 & -0.00041666 & 0.00269816 \\
3;1 & W & 6 & -0.000850866 & 0.00225551 \\
3;1 & W & 7 & -0.0023459 & 0.000745795 \\
3;1 & W & 8 & -0.00124523 & 0.00180352 \\
3;1 & W & 9 & -0.00147051 & 0.00156613 \\
3;1 & W & 10 & -0.00150886 & 0.00151075 \\
3;1 & W & 11 & -0.00150001 & 0.00150822 \\
3;1 & W & 12 & -0.00150165 & 0.00149869 \\
3;1 & W & 13 & -0.00149858 & 0.00149581 \\
3;1 & W & 14 & -0.00149663 & 0.00149349 \\
3;1 & W & 15 & -0.00149339 & 0.00149403 \\
3;1 & W & 16 & 0 & 0 \\
3;1 & F & 4 & 0.000691523 & 0.00372718 \\
3;1 & F & 5 & -0.000505628 & 0.00250065 \\
3;1 & F & 6 & -0.0010836 & 0.00190806 \\
3;1 & F & 7 & -0.00122457 & 0.00176312 \\
3;1 & F & 8 & -0.001473 & 0.00150813 \\
3;1 & F & 9 & -0.0014817 & 0.00149918 \\
3;1 & F & 10 & -0.0014861 & 0.00149463 \\
3;1 & F & 11 & -0.00148819 & 0.00149247 \\
3;1 & F & 12 & -0.00148929 & 0.00149132 \\
3;1 & F & 13 & -0.00148981 & 0.00149079 \\
3;1 & F & 14 & -0.00149008 & 0.0014905 \\
3;1 & F & 15 & -0.00149022 & 0.00149036 \\
3;1 & F & 16 & 0 & 0 \\
3;4 & W & 4 & -0.00115463 & 0.00216779 \\
3;4 & W & 5 & -0.00128854 & 0.00202944 \\
3;4 & W & 6 & -0.0012672 & 0.00204406 \\
3;4 & W & 7 & -0.00203778 & 0.00126435 \\
3;4 & W & 8 & -0.00147323 & 0.00179482 \\
3;4 & W & 9 & -0.00162558 & 0.00163607 \\
3;4 & W & 10 & -0.00162825 & 0.00162325 \\
3;4 & W & 11 & -0.00162181 & 0.00162209 \\
3;4 & W & 12 & -0.00161759 & 0.00162064 \\
3;4 & W & 13 & -0.00161723 & 0.00161692 \\
3;4 & W & 14 & -0.00161637 & 0.00161497 \\
3;4 & W & 15 & -0.00161356 & 0.00161547 \\
3;4 & W & 16 & 0 & 0 \\
3;4 & F & 4 & -0.00143754 & 0.00179918 \\
3;4 & F & 5 & -0.0015455 & 0.00168516 \\
3;4 & F & 6 & -0.0015881 & 0.00163939 \\
3;4 & F & 7 & -0.00159646 & 0.0016305 \\
3;4 & F & 8 & -0.00161162 & 0.00161347 \\
3;4 & F & 9 & -0.00161206 & 0.00161302 \\
3;4 & F & 10 & -0.00161229 & 0.00161278 \\
3;4 & F & 11 & -0.00161241 & 0.00161265 \\
3;4 & F & 12 & -0.00161247 & 0.00161258 \\
3;4 & F & 13 & -0.0016125 & 0.00161255 \\
3;4 & F & 14 & -0.00161251 & 0.00161254 \\
3;4 & F & 15 & -0.00161252 & 0.00161253 \\
3;4 & F & 16 & 0 & 0 \\
\end{longtable}
}
\noindent{\footnotesize All rows use the uniform initial law and the actual first-target-one controller. Self-comparisons are exact identities. Both policies use the declared robust repair and action quantum.\par}


\section{Reconstruction and preservation}\label{supp:audit49}
The release verifies frozen science-source identities, every recorded checkpoint, all signed policy formulas, every first crossing, complete-prefix counts, repeated policy identities and the confidence intervals reconstructed from endpoint moments. Additional exact tests check the primitive regions, resource allocation, original-owner identity, fee semantics and net-choice bound. The complete development article and supplement retain all baseline labeled results; the original paths and adverse results are unchanged. A clean Git-archive rebuild removes generated PDFs and derived tables, then rebuilds from ordinary committed sources and frozen records without network or new scientific execution. Compilation, finite tests and source identities do not constitute independent mathematical peer approval.
\section{Every R50 actual-cost interval and mechanism record}\label{supp:records50}
The following tables use the corrected outward-subtraction replay. C denotes the separate common-accuracy block, P the primary target-five plan, Q the primary target-two plan, O the original R49 target-two policies, A residual-driven FVI, and I the isotropic block. U is the uniform initial law; L, M and H denote point laws at 1/8,1/2 and 7/8 in every coordinate. Each group has four absolute costs and six contrasts; the table lists all 280 estimands. Printed endpoints are rounded for display. Exact rational endpoints and source-bound moments govern inference.
{\small
\setlength{\tabcolsep}{3pt}
\begin{longtable}{@{}llrr@{}}
\caption{Every new absolute expected policy-cost interval}\label{tab:all-costs50}\\
\toprule
Case & Policy & Lower & Upper \\
\midrule
\endfirsthead
\toprule
Case & Policy & Lower & Upper \\
\midrule
\endhead
\bottomrule
\endfoot
C2;2;1;U & left & 0.6242 & 0.6378 \\
C2;2;1;U & left-repaired & 0.61117 & 0.62482 \\
C2;2;1;U & right & 0.59821 & 0.61177 \\
C2;2;1;U & right-repaired & 0.59628 & 0.60988 \\
C3;2;1;U & left & 0.58992 & 0.60068 \\
C3;2;1;U & left-repaired & 0.58062 & 0.59146 \\
C3;2;1;U & right & 0.56904 & 0.57982 \\
C3;2;1;U & right-repaired & 0.56768 & 0.57849 \\
C4;2;1;U & left & 0.62188 & 0.63199 \\
C4;2;1;U & left-repaired & 0.60399 & 0.61391 \\
C4;2;1;U & right & 0.5855 & 0.59525 \\
C4;2;1;U & right-repaired & 0.5854 & 0.59515 \\
I2;2;1;U & left & 0.63959 & 0.65386 \\
I2;2;1;U & left-repaired & 0.61456 & 0.6285 \\
I2;2;1;U & right & 0.59407 & 0.60781 \\
I2;2;1;U & right-repaired & 0.59363 & 0.60737 \\
O2;2;1;M & left & 0.30128 & 0.30382 \\
O2;2;1;M & left-repaired & 0.29236 & 0.29488 \\
O2;2;1;M & right & 0.29238 & 0.29489 \\
O2;2;1;M & right-repaired & 0.29236 & 0.29488 \\
O2;2;1;L & left & 1.5737 & 1.5766 \\
O2;2;1;L & left-repaired & 1.5626 & 1.565 \\
O2;2;1;L & right & 1.5163 & 1.5188 \\
O2;2;1;L & right-repaired & 1.5162 & 1.5186 \\
O2;2;1;H & left & 0.19157 & 0.19409 \\
O2;2;1;H & left-repaired & 0.18888 & 0.19135 \\
O2;2;1;H & right & 0.18497 & 0.18741 \\
O2;2;1;H & right-repaired & 0.18423 & 0.18668 \\
O2;2;1;U & left & 0.61089 & 0.62464 \\
O2;2;1;U & left-repaired & 0.60392 & 0.61765 \\
O2;2;1;U & right & 0.59513 & 0.60882 \\
O2;2;1;U & right-repaired & 0.59491 & 0.6086 \\
O2;2;4;M & left & 0.42746 & 0.43017 \\
O2;2;4;M & left-repaired & 0.42249 & 0.42507 \\
O2;2;4;M & right & 0.38962 & 0.39224 \\
O2;2;4;M & right-repaired & 0.38952 & 0.39214 \\
O2;2;4;L & left & 1.6588 & 1.6614 \\
O2;2;4;L & left-repaired & 1.6564 & 1.659 \\
O2;2;4;L & right & 1.6358 & 1.6385 \\
O2;2;4;L & right-repaired & 1.6358 & 1.6385 \\
O2;2;4;H & left & 0.2474 & 0.2501 \\
O2;2;4;H & left-repaired & 0.24264 & 0.24526 \\
O2;2;4;H & right & 0.20559 & 0.20818 \\
O2;2;4;H & right-repaired & 0.20554 & 0.20813 \\
O2;2;4;U & left & 0.69261 & 0.7071 \\
O2;2;4;U & left-repaired & 0.68704 & 0.70155 \\
O2;2;4;U & right & 0.67867 & 0.69319 \\
O2;2;4;U & right-repaired & 0.67849 & 0.69301 \\
O2;3;1;M & left & 0.42649 & 0.42982 \\
O2;3;1;M & left-repaired & 0.42326 & 0.42657 \\
O2;3;1;M & right & 0.40778 & 0.41096 \\
O2;3;1;M & right-repaired & 0.40776 & 0.41095 \\
O2;3;1;L & left & 1.6968 & 1.7 \\
O2;3;1;L & left-repaired & 1.6932 & 1.6964 \\
O2;3;1;L & right & 1.6673 & 1.6705 \\
O2;3;1;L & right-repaired & 1.6672 & 1.6704 \\
O2;3;1;H & left & 0.29136 & 0.29449 \\
O2;3;1;H & left-repaired & 0.28954 & 0.29265 \\
O2;3;1;H & right & 0.27712 & 0.28025 \\
O2;3;1;H & right-repaired & 0.27696 & 0.28009 \\
O2;3;1;U & left & 0.71942 & 0.73411 \\
O2;3;1;U & left-repaired & 0.71638 & 0.73106 \\
O2;3;1;U & right & 0.70681 & 0.72143 \\
O2;3;1;U & right-repaired & 0.70676 & 0.72139 \\
O2;3;4;M & left & 0.58099 & 0.58436 \\
O2;3;4;M & left-repaired & 0.57925 & 0.58262 \\
O2;3;4;M & right & 0.56732 & 0.57065 \\
O2;3;4;M & right-repaired & 0.56709 & 0.57042 \\
O2;3;4;L & left & 1.8582 & 1.8616 \\
O2;3;4;L & left-repaired & 1.856 & 1.8594 \\
O2;3;4;L & right & 1.8454 & 1.8487 \\
O2;3;4;L & right-repaired & 1.8454 & 1.8487 \\
O2;3;4;H & left & 0.35451 & 0.35788 \\
O2;3;4;H & left-repaired & 0.35176 & 0.35509 \\
O2;3;4;H & right & 0.34392 & 0.34723 \\
O2;3;4;H & right-repaired & 0.34391 & 0.34721 \\
O2;3;4;U & left & 0.86235 & 0.87799 \\
O2;3;4;U & left-repaired & 0.86048 & 0.87613 \\
O2;3;4;U & right & 0.85339 & 0.86906 \\
O2;3;4;U & right-repaired & 0.85323 & 0.86891 \\
Q2;2;1;U & left & 0.60972 & 0.62332 \\
Q2;2;1;U & left-repaired & 0.60276 & 0.61638 \\
Q2;2;1;U & right & 0.59518 & 0.60882 \\
Q2;2;1;U & right-repaired & 0.59444 & 0.6081 \\
P2;2;1;U & left & 0.6222 & 0.63577 \\
P2;2;1;U & left-repaired & 0.60912 & 0.62274 \\
P2;2;1;U & right & 0.59639 & 0.60993 \\
P2;2;1;U & right-repaired & 0.59445 & 0.60803 \\
Q2;2;4;U & left & 0.69243 & 0.70702 \\
Q2;2;4;U & left-repaired & 0.68754 & 0.70213 \\
Q2;2;4;U & right & 0.68106 & 0.69562 \\
Q2;2;4;U & right-repaired & 0.68086 & 0.69542 \\
Q2;3;1;U & left & 0.72237 & 0.73696 \\
Q2;3;1;U & left-repaired & 0.71906 & 0.73365 \\
Q2;3;1;U & right & 0.70782 & 0.7224 \\
Q2;3;1;U & right-repaired & 0.70769 & 0.72228 \\
Q2;3;4;U & left & 0.86908 & 0.88468 \\
Q2;3;4;U & left-repaired & 0.86539 & 0.881 \\
Q2;3;4;U & right & 0.85511 & 0.87074 \\
Q2;3;4;U & right-repaired & 0.85451 & 0.87015 \\
P3;2;1;U & left & 0.58981 & 0.60055 \\
P3;2;1;U & left-repaired & 0.58046 & 0.59128 \\
P3;2;1;U & right & 0.56902 & 0.57978 \\
P3;2;1;U & right-repaired & 0.56764 & 0.57844 \\
P4;2;1;U & left & 0.61951 & 0.62957 \\
P4;2;1;U & left-repaired & 0.60168 & 0.61155 \\
P4;2;1;U & right & 0.58357 & 0.59328 \\
P4;2;1;U & right-repaired & 0.58332 & 0.59303 \\
A2;2;1;U & left & 0.59772 & 0.61129 \\
A2;2;1;U & left-repaired & 0.59585 & 0.60945 \\
A2;2;1;U & right & 0.59777 & 0.61134 \\
A2;2;1;U & right-repaired & 0.59585 & 0.60946 \\
\end{longtable}
}

{\small
\setlength{\tabcolsep}{3pt}
\begin{longtable}{@{}llrr@{}}
\caption{Every new direct expected policy-cost contrast}\label{tab:all-contrasts50}\\
\toprule
Case & Contrast & Lower & Upper \\
\midrule
\endfirsthead
\toprule
Case & Contrast & Lower & Upper \\
\midrule
\endhead
\bottomrule
\endfoot
C2;2;1;U & L minus L+ & 0.011448 & 0.014567 \\
C2;2;1;U & L minus R & 0.023033 & 0.028988 \\
C2;2;1;U & L minus R+ & 0.024948 & 0.030894 \\
C2;2;1;U & L+ minus R & 0.010133 & 0.015872 \\
C2;2;1;U & L+ minus R+ & 0.012051 & 0.017775 \\
C2;2;1;U & R minus R+ & 0.00067113 & 0.0031495 \\
C3;2;1;U & L minus L+ & 0.0077535 & 0.010757 \\
C3;2;1;U & L minus R & 0.01796 & 0.023776 \\
C3;2;1;U & L minus R+ & 0.019309 & 0.025125 \\
C3;2;1;U & L+ minus R & 0.0087758 & 0.014449 \\
C3;2;1;U & L+ minus R+ & 0.010129 & 0.015794 \\
C3;2;1;U & R minus R+ & 0.0001224 & 0.0025755 \\
C4;2;1;U & L minus L+ & 0.016432 & 0.019542 \\
C4;2;1;U & L minus R & 0.033636 & 0.039473 \\
C4;2;1;U & L minus R+ & 0.033741 & 0.039578 \\
C4;2;1;U & L+ minus R & 0.015807 & 0.021327 \\
C4;2;1;U & L+ minus R+ & 0.015912 & 0.021433 \\
C4;2;1;U & R minus R+ & 0 & 0.0012997 \\
I2;2;1;U & L minus L+ & 0.023541 & 0.026849 \\
I2;2;1;U & L minus R & 0.042713 & 0.048868 \\
I2;2;1;U & L minus R+ & 0.043151 & 0.049305 \\
I2;2;1;U & L+ minus R & 0.017772 & 0.023418 \\
I2;2;1;U & L+ minus R+ & 0.018209 & 0.023856 \\
I2;2;1;U & R minus R+ & 0 & 0.00164 \\
O2;2;1;M & L minus L+ & 0.007622 & 0.010239 \\
O2;2;1;M & L minus R & 0.0064146 & 0.011415 \\
O2;2;1;M & L minus R+ & 0.00643 & 0.011431 \\
O2;2;1;M & L+ minus R & -0.0023997 & 0.002369 \\
O2;2;1;M & L+ minus R+ & -0.0023843 & 0.0023843 \\
O2;2;1;M & R minus R+ & 0 & 0.0012077 \\
O2;2;1;L & L minus L+ & 0.0099411 & 0.012722 \\
O2;2;1;L & L minus R & 0.054972 & 0.060123 \\
O2;2;1;L & L minus R+ & 0.055155 & 0.060307 \\
O2;2;1;L & L+ minus R & 0.043822 & 0.048609 \\
O2;2;1;L & L+ minus R+ & 0.044007 & 0.048792 \\
O2;2;1;L & R minus R+ & 0 & 0.001381 \\
O2;2;1;H & L minus L+ & 0.0014934 & 0.003946 \\
O2;2;1;H & L minus R & 0.0042173 & 0.0090667 \\
O2;2;1;H & L minus R+ & 0.0049502 & 0.0097948 \\
O2;2;1;H & L+ minus R & 0.001523 & 0.0063216 \\
O2;2;1;H & L+ minus R+ & 0.0022567 & 0.007049 \\
O2;2;1;H & R minus R+ & 0 & 0.0019272 \\
O2;2;1;U & L minus L+ & 0.0056474 & 0.0083165 \\
O2;2;1;U & L minus R & 0.013161 & 0.018416 \\
O2;2;1;U & L minus R+ & 0.013385 & 0.01864 \\
O2;2;1;U & L+ minus R & 0.0062318 & 0.011381 \\
O2;2;1;U & L+ minus R+ & 0.0064562 & 0.011605 \\
O2;2;1;U & R minus R+ & 0 & 0.0014206 \\
O2;2;4;M & L minus L+ & 0.0036933 & 0.0063836 \\
O2;2;4;M & L minus R & 0.035287 & 0.040496 \\
O2;2;4;M & L minus R+ & 0.035387 & 0.040597 \\
O2;2;4;M & L+ minus R & 0.03028 & 0.035426 \\
O2;2;4;M & L+ minus R+ & 0.030381 & 0.035526 \\
O2;2;4;M & R minus R+ & 0 & 0.0013773 \\
O2;2;4;L & L minus L+ & 0.0010828 & 0.0036974 \\
O2;2;4;L & L minus R & 0.020345 & 0.025506 \\
O2;2;4;L & L minus R+ & 0.020345 & 0.025506 \\
O2;2;4;L & L+ minus R & 0.017976 & 0.023095 \\
O2;2;4;L & L+ minus R+ & 0.017976 & 0.023095 \\
O2;2;4;L & R minus R+ & 0 & 0.0012764 \\
O2;2;4;H & L minus L+ & 0.0034416 & 0.006155 \\
O2;2;4;H & L minus R & 0.039236 & 0.044497 \\
O2;2;4;H & L minus R+ & 0.039288 & 0.044549 \\
O2;2;4;H & L+ minus R & 0.034497 & 0.039639 \\
O2;2;4;H & L+ minus R+ & 0.034549 & 0.039692 \\
O2;2;4;H & R minus R+ & 0 & 0.0013293 \\
O2;2;4;U & L minus L+ & 0.004188 & 0.0069254 \\
O2;2;4;U & L minus R & 0.011181 & 0.016658 \\
O2;2;4;U & L minus R+ & 0.011361 & 0.016838 \\
O2;2;4;U & L+ minus R & 0.0056495 & 0.011076 \\
O2;2;4;U & L+ minus R+ & 0.0058295 & 0.011256 \\
O2;2;4;U & R minus R+ & 0 & 0.0014585 \\
O2;3;1;M & L minus L+ & 0.0016956 & 0.0047777 \\
O2;3;1;M & L minus R & 0.015697 & 0.021868 \\
O2;3;1;M & L minus R+ & 0.015714 & 0.021886 \\
O2;3;1;M & L+ minus R & 0.012467 & 0.018624 \\
O2;3;1;M & L+ minus R+ & 0.012485 & 0.018641 \\
O2;3;1;M & R minus R+ & 0 & 0.0015083 \\
O2;3;1;L & L minus L+ & 0.0020031 & 0.005116 \\
O2;3;1;L & L minus R & 0.026367 & 0.032516 \\
O2;3;1;L & L minus R+ & 0.026477 & 0.032626 \\
O2;3;1;L & L+ minus R & 0.022819 & 0.028945 \\
O2;3;1;L & L+ minus R+ & 0.022928 & 0.029055 \\
O2;3;1;L & R minus R+ & 0 & 0.0016012 \\
O2;3;1;H & L minus L+ & 0.00030976 & 0.0033545 \\
O2;3;1;H & L minus R & 0.011221 & 0.017268 \\
O2;3;1;H & L minus R+ & 0.011378 & 0.017424 \\
O2;3;1;H & L+ minus R & 0.0093942 & 0.01543 \\
O2;3;1;H & L+ minus R+ & 0.0095515 & 0.015586 \\
O2;3;1;H & R minus R+ & 0 & 0.0016481 \\
O2;3;1;U & L minus L+ & 0.0014967 & 0.0045926 \\
O2;3;1;U & L minus R & 0.0095127 & 0.015773 \\
O2;3;1;U & L minus R+ & 0.0095603 & 0.01582 \\
O2;3;1;U & L+ minus R & 0.0064826 & 0.012713 \\
O2;3;1;U & L+ minus R+ & 0.0065302 & 0.012761 \\
O2;3;1;U & R minus R+ & 0 & 0.0015385 \\
O2;3;4;M & L minus L+ & 9.4671e-05 & 0.0033779 \\
O2;3;4;M & L minus R & 0.010433 & 0.016948 \\
O2;3;4;M & L minus R+ & 0.01066 & 0.017175 \\
O2;3;4;M & L+ minus R & 0.0087064 & 0.015202 \\
O2;3;4;M & L+ minus R+ & 0.0089337 & 0.015429 \\
O2;3;4;M & R minus R+ & 0 & 0.00184 \\
O2;3;4;L & L minus L+ & 0.0005048 & 0.0037993 \\
O2;3;4;L & L minus R & 0.0095446 & 0.01608 \\
O2;3;4;L & L minus R+ & 0.0095504 & 0.016086 \\
O2;3;4;L & L+ minus R & 0.0074107 & 0.01391 \\
O2;3;4;L & L+ minus R+ & 0.0074164 & 0.013916 \\
O2;3;4;L & R minus R+ & 0 & 0.0016183 \\
O2;3;4;H & L minus L+ & 0.0011154 & 0.0044219 \\
O2;3;4;H & L minus R & 0.0073415 & 0.013898 \\
O2;3;4;H & L minus R+ & 0.0073538 & 0.013911 \\
O2;3;4;H & L+ minus R & 0.004596 & 0.011107 \\
O2;3;4;H & L+ minus R+ & 0.0046083 & 0.011119 \\
O2;3;4;H & R minus R+ & 0 & 0.0016249 \\
O2;3;4;U & L minus L+ & 0.00021625 & 0.0035063 \\
O2;3;4;U & L minus R & 0.0056055 & 0.012273 \\
O2;3;4;U & L minus R+ & 0.0057651 & 0.012433 \\
O2;3;4;U & L+ minus R & 0.0037509 & 0.010406 \\
O2;3;4;U & L+ minus R+ & 0.0039105 & 0.010565 \\
O2;3;4;U & R minus R+ & 0 & 0.001774 \\
Q2;2;1;U & L minus L+ & 0.0055893 & 0.0083122 \\
Q2;2;1;U & L minus R & 0.01186 & 0.017187 \\
Q2;2;1;U & L minus R+ & 0.01259 & 0.017917 \\
Q2;2;1;U & L+ minus R & 0.0049684 & 0.010177 \\
Q2;2;1;U & L+ minus R+ & 0.0056983 & 0.010907 \\
Q2;2;1;U & R minus R+ & 0 & 0.0019382 \\
P2;2;1;U & L minus L+ & 0.011497 & 0.014615 \\
P2;2;1;U & L minus R & 0.022852 & 0.028801 \\
P2;2;1;U & L minus R+ & 0.024778 & 0.030719 \\
P2;2;1;U & L+ minus R & 0.0099052 & 0.015635 \\
P2;2;1;U & L+ minus R+ & 0.011835 & 0.01755 \\
P2;2;1;U & R minus R+ & 0.00068331 & 0.0031614 \\
Q2;2;4;U & L minus L+ & 0.003535 & 0.0062382 \\
Q2;2;4;U & L minus R & 0.0086864 & 0.014074 \\
Q2;2;4;U & L minus R+ & 0.0088869 & 0.014274 \\
Q2;2;4;U & L+ minus R & 0.0038216 & 0.0091655 \\
Q2;2;4;U & L+ minus R+ & 0.004022 & 0.0093656 \\
Q2;2;4;U & R minus R+ & 0 & 0.0014794 \\
Q2;3;1;U & L minus L+ & 0.0017403 & 0.0048791 \\
Q2;3;1;U & L minus R & 0.01136 & 0.017753 \\
Q2;3;1;U & L minus R+ & 0.011488 & 0.01788 \\
Q2;3;1;U & L+ minus R & 0.0080678 & 0.014426 \\
Q2;3;1;U & L+ minus R+ & 0.0081953 & 0.014553 \\
Q2;3;1;U & R minus R+ & 0 & 0.0016207 \\
Q2;3;4;U & L minus L+ & 0.0020091 & 0.0053568 \\
Q2;3;4;U & L minus R & 0.010559 & 0.017349 \\
Q2;3;4;U & L minus R+ & 0.011154 & 0.017943 \\
Q2;3;4;U & L+ minus R & 0.006888 & 0.013654 \\
Q2;3;4;U & L+ minus R+ & 0.0074835 & 0.014248 \\
Q2;3;4;U & R minus R+ & 0 & 0.0022148 \\
P3;2;1;U & L minus L+ & 0.0078066 & 0.010814 \\
P3;2;1;U & L minus R & 0.017874 & 0.023687 \\
P3;2;1;U & L minus R+ & 0.019232 & 0.025044 \\
P3;2;1;U & L+ minus R & 0.008637 & 0.014304 \\
P3;2;1;U & L+ minus R+ & 0.0099988 & 0.015657 \\
P3;2;1;U & R minus R+ & 0.00013073 & 0.0025845 \\
P4;2;1;U & L minus L+ & 0.016374 & 0.019482 \\
P4;2;1;U & L minus R & 0.033196 & 0.039029 \\
P4;2;1;U & L minus R+ & 0.033446 & 0.039279 \\
P4;2;1;U & L+ minus R & 0.015426 & 0.020943 \\
P4;2;1;U & L+ minus R+ & 0.015676 & 0.021193 \\
P4;2;1;U & R minus R+ & 0 & 0.0014477 \\
A2;2;1;U & L minus L+ & 0.00062241 & 0.0030972 \\
A2;2;1;U & L minus R & -0.0024517 & 0.0023493 \\
A2;2;1;U & L minus R+ & -0.00057956 & 0.0042851 \\
A2;2;1;U & L+ minus R & -0.0043435 & 0.00052149 \\
A2;2;1;U & L+ minus R+ & -0.0024125 & 0.0023985 \\
A2;2;1;U & R minus R+ & 0.00066516 & 0.0031428 \\
\end{longtable}
}
\noindent{\footnotesize Both tables are reconstructed from stored outward means, variance bounds and deterministic supports. Numerical replay uses the same original samples. Statistical interpretation requires the independent-bin model.\par}

{\small
\setlength{\tabcolsep}{3pt}
\begin{longtable}{@{}llrrrrr@{}}
\caption{Final-action mechanism diagnostics}\label{tab:mechanism50}\\
\toprule
Case & Policy & Accepted & Blocked & Old action & New action & Derivative calls \\
\midrule
\endfirsthead
\toprule
Case & Policy & Accepted & Blocked & Old action & New action & Derivative calls \\
\midrule
\endhead
\bottomrule
\endfoot
C2;2;1;U & L & 81251 & 1686 & 0.13805 & 0.16199 & 1216761 \\
C2;2;1;U & R & 92005 & 2261 & 0.17813 & 0.16288 & 1223163 \\
C3;2;1;U & L & 71835 & 1430 & 0.15264 & 0.16793 & 1212973 \\
C3;2;1;U & R & 86739 & 1858 & 0.18142 & 0.16825 & 1220030 \\
C4;2;1;U & L & 105827 & 1124 & 0.11931 & 0.16893 & 1209765 \\
C4;2;1;U & R & 79985 & 9020 & 0.16899 & 0.16935 & 1219031 \\
I2;2;1;U & L & 115800 & 1824 & 0.10471 & 0.16094 & 1216727 \\
I2;2;1;U & R & 87457 & 6803 & 0.16155 & 0.16239 & 1222172 \\
O2;2;1;M & L & 123097 & 0 & 0.13434 & 0.18236 & 1179648 \\
O2;2;1;M & R & 66240 & 0 & 0.18188 & 0.18236 & 1179648 \\
O2;2;1;L & L & 51516 & 0 & 0.13357 & 0.14812 & 1179648 \\
O2;2;1;L & R & 25162 & 0 & 0.14947 & 0.14977 & 1179648 \\
O2;2;1;H & L & 131072 & 0 & 0.11338 & 0.11033 & 1301723 \\
O2;2;1;H & R & 131072 & 0 & 0.10626 & 0.12107 & 1286030 \\
O2;2;1;U & L & 96199 & 2180 & 0.13976 & 0.16238 & 1217297 \\
O2;2;1;U & R & 86179 & 5606 & 0.16295 & 0.16247 & 1223021 \\
O2;2;4;M & L & 131072 & 0 & 0.06499 & 0.09143 & 1253562 \\
O2;2;4;M & R & 131072 & 0 & 0.08886 & 0.08423 & 1245289 \\
O2;2;4;L & L & 131072 & 0 & 0.12539 & 0.13479 & 1227082 \\
O2;2;4;L & R & 0 & 131072 & 0.13106 & 0.13106 & 1232486 \\
O2;2;4;H & L & 131072 & 0 & 0.04179 & 0.04061 & 1294458 \\
O2;2;4;H & R & 131072 & 0 & 0.05060 & 0.04725 & 1278365 \\
O2;2;4;U & L & 127091 & 2453 & 0.08431 & 0.08510 & 1255705 \\
O2;2;4;U & R & 116259 & 12245 & 0.08476 & 0.08500 & 1255578 \\
O2;3;1;M & L & 74044 & 0 & 0.15964 & 0.17640 & 1179648 \\
O2;3;1;M & R & 50391 & 0 & 0.17885 & 0.17907 & 1179648 \\
O2;3;1;L & L & 47744 & 0 & 0.14872 & 0.15570 & 1179648 \\
O2;3;1;L & R & 68952 & 0 & 0.15623 & 0.15647 & 1179648 \\
O2;3;1;H & L & 127923 & 3149 & 0.15631 & 0.16935 & 1267556 \\
O2;3;1;H & R & 130156 & 880 & 0.16622 & 0.15981 & 1272977 \\
O2;3;1;U & L & 93095 & 1636 & 0.15635 & 0.17217 & 1200960 \\
O2;3;1;U & R & 68286 & 6835 & 0.17297 & 0.17285 & 1207373 \\
O2;3;4;M & L & 131009 & 63 & 0.08959 & 0.08991 & 1253921 \\
O2;3;4;M & R & 131072 & 0 & 0.08589 & 0.09307 & 1249060 \\
O2;3;4;L & L & 129450 & 1622 & 0.11835 & 0.11804 & 1234478 \\
O2;3;4;L & R & 75986 & 44091 & 0.11664 & 0.11708 & 1237417 \\
O2;3;4;H & L & 130535 & 487 & 0.06486 & 0.07197 & 1266116 \\
O2;3;4;H & R & 81893 & 46520 & 0.06993 & 0.07110 & 1270467 \\
O2;3;4;U & L & 127548 & 3104 & 0.09220 & 0.09283 & 1251214 \\
O2;3;4;U & R & 119926 & 9604 & 0.09243 & 0.09264 & 1251454 \\
Q2;2;1;U & L & 84626 & 1862 & 0.14266 & 0.16266 & 1217576 \\
Q2;2;1;U & R & 89341 & 2780 & 0.16443 & 0.16248 & 1222725 \\
P2;2;1;U & L & 81414 & 1734 & 0.13768 & 0.16190 & 1217185 \\
P2;2;1;U & R & 92131 & 2223 & 0.17824 & 0.16278 & 1223408 \\
Q2;2;4;U & L & 127314 & 3044 & 0.08473 & 0.08526 & 1255543 \\
Q2;2;4;U & R & 116513 & 12198 & 0.08462 & 0.08509 & 1255620 \\
Q2;3;1;U & L & 78872 & 1217 & 0.15697 & 0.17211 & 1200463 \\
Q2;3;1;U & R & 72138 & 3059 & 0.17400 & 0.17287 & 1207783 \\
Q2;3;4;U & L & 127944 & 2619 & 0.09086 & 0.09284 & 1251053 \\
Q2;3;4;U & R & 124704 & 5413 & 0.09262 & 0.09261 & 1251454 \\
P3;2;1;U & L & 71967 & 1425 & 0.15260 & 0.16790 & 1213011 \\
P3;2;1;U & R & 86489 & 1881 & 0.18143 & 0.16824 & 1220010 \\
P4;2;1;U & L & 106376 & 1130 & 0.11948 & 0.16893 & 1210149 \\
P4;2;1;U & R & 86994 & 7077 & 0.16745 & 0.16930 & 1218786 \\
A2;2;1;U & L & 90517 & 2162 & 0.17906 & 0.16290 & 1223395 \\
A2;2;1;U & R & 92064 & 2277 & 0.17821 & 0.16286 & 1223276 \\
\end{longtable}
}
\noindent{\footnotesize Means are descriptive over the recorded paths, not confidence intervals for separate mechanism parameters. Complete date-wise earlier actions, terminal state vectors, capacity and quantum repair counts, and exact fallbacks are deposited.\par}

{\small
\setlength{\tabcolsep}{3pt}
\begin{longtable}{@{}llrrrrr@{}}
\caption{New-construction operation and storage account}\label{tab:work50}\\
\toprule
Case & Method & Target eval. & Actor scalars & KiB & Coeff. bits & MiB \\
\midrule
\endfirsthead
\toprule
Case & Method & Target eval. & Actor scalars & KiB & Coeff. bits & MiB \\
\midrule
\endhead
\bottomrule
\endfoot
C2;2;1;U & W & 1,156 & 578 & 41.8 & 36 & 41.4 \\
C2;2;1;U & F & 1,156 & 578 & 26.1 & 42 & 41.7 \\
C3;2;1;U & W & 19,652 & 9,826 & 851.1 & 42 & 44.3 \\
C3;2;1;U & F & 19,652 & 9,826 & 434.4 & 42 & 42.3 \\
C4;2;1;U & W & 236,196 & 13,122 & 1024.5 & 42 & 50.5 \\
C4;2;1;U & F & 3,242,952 & 167,042 & 7464.3 & 42 & 169.8 \\
I2;2;1;U & W & 11,664 & 162 & 14.2 & 41 & 42.4 \\
I2;2;1;U & F & 11,664 & 162 & 9.3 & 41 & 42.0 \\
Q2;2;1;U & W & 13,068 & 2,178 & 171.9 & 41 & 41.6 \\
Q2;2;1;U & F & 13,068 & 2,178 & 98.6 & 42 & 42.1 \\
P2;2;1;U & W & 1,156 & 578 & 41.8 & 36 & 42.0 \\
P2;2;1;U & F & 1,156 & 578 & 26.1 & 42 & 42.0 \\
Q2;2;4;U & W & 19,602 & 2,178 & 178.4 & 42 & 41.6 \\
Q2;2;4;U & F & 19,602 & 2,178 & 99.4 & 42 & 42.0 \\
Q2;3;1;U & W & 126,750 & 12,675 & 1005.0 & 42 & 45.8 \\
Q2;3;1;U & F & 126,750 & 12,675 & 550.0 & 42 & 41.9 \\
Q2;3;4;U & W & 126,750 & 12,675 & 983.8 & 42 & 45.8 \\
Q2;3;4;U & F & 126,750 & 12,675 & 539.4 & 42 & 41.8 \\
P3;2;1;U & W & 19,652 & 9,826 & 851.1 & 42 & 44.5 \\
P3;2;1;U & F & 19,652 & 9,826 & 434.4 & 42 & 42.5 \\
P4;2;1;U & W & 236,196 & 13,122 & 1024.5 & 42 & 51.3 \\
P4;2;1;U & F & 236,196 & 13,122 & 575.2 & 42 & 48.3 \\
A2;2;1;U & A & 4,472 & 578 & 39.5 & 42 & 41.9 \\
A2;2;1;U & F & 1,156 & 578 & 26.1 & 42 & 42.0 \\
\end{longtable}
}
\noindent{\footnotesize Target evaluations include every failed common-accuracy attempt. Actor scalars and checkpoint KiB describe the returned model. Coefficient bits are the true maximum, not the additive descriptor in the raw prefix map; MiB is process peak memory. Full operation categories, including terminal gates and exact fallbacks, remain in JSON.\par}


\section{Numerical correction and sequential common-accuracy design}\label{supp:corrections50}
The primary shared allocation is sufficient for the witness bound, not automatically for conventional FVI's actual interpolated modulus. Its failed certificates are retained. The common-accuracy amendment was specified after that outcome and frozen before its own execution; it charges every failed construction before doubling the state quota. It is not a retrospective recoding of the primary study.

During publication review, nearest-rounded subtraction of different policies' already outward path endpoints was replaced by directed outward subtraction. The same bin streams, policies and sample counts were replayed. Original sources and results remain unchanged; a separate correction manifest binds the new evaluator, and the audit checks identical streams and unchanged sign classifications. Replaying the same observations adds no statistical sample size. The additive raw prefix field for coefficient-bit descriptors is not interpreted as a maximum; maximum bit length and live storage are reconstructed separately. The correction and development logs are part of the release.

\section{Exhaustive robustness-service accounting}\label{supp:work52}
\begin{table}[htbp]
\caption{Complete robustness-service work and memory}
\label{tab:work52}
\centering\small
\begin{tabular}{lrrrr}
\toprule
Incumbent & $T$ & $p$ & Seconds & Peak MiB \\
\midrule
Witness & 2 & 1 & 3.73 & 50.6 \\
Witness & 2 & 4 & 3.72 & 50.6 \\
Witness & 3 & 1 & 4.37 & 56.5 \\
Witness & 3 & 4 & 4.32 & 56.5 \\
FVI & 2 & 1 & 3.63 & 50.7 \\
FVI & 2 & 4 & 3.76 & 51.0 \\
FVI & 3 & 1 & 3.91 & 52.5 \\
FVI & 3 & 4 & 3.86 & 52.3 \\
\bottomrule
\end{tabular}
\parbox{0.96\textwidth}{\footnotesize Each isolated service loads a frozen incumbent, computes all proposals and centered intervals, independently replays the inherited terminal gate, evaluates every amplitude gate, and durably serializes the cell records. Process startup is excluded; inherited construction time is not added. These are verification-service clocks, not a work-to-policy-cost superiority claim.}
\end{table}

\section{Directed full-sweep verification records}\label{supp:directed54}
The new main theorem is proved in the article. This supplement records the numerical objects to which it is applied. The scientific implementation and both source freezes are retained from the R53 execution. The primary services reconstruct the original two-state checkpoints from primitive costs, transitions, action menus and integration rules before building acquired actors. All subsequent policies are integer arrays of downward action indices. Their complete byte identities, rather than only their expected costs, determine whether a later pass has changed a policy.

For every date and pass, the raw cell archive contains nine candidate lower and upper advantage endpoints, the candidate quantum indices, eight lower-cover intervals, the selected action, the accepted upper endpoint, the lower bound on the continuous-action infimum and the directed greedy-gap upper endpoint. The publication replay reconstructs the strict incumbent-preserving selection and compares it with the deployed array. It also recomputes the maximum directed gap and the exact rational finite-sweep recursion. Feasibility is checked using the same complete-cell lower-capacity integer formula, independently of a displayed decimal capacity.

\begin{table}[htbp]
\centering
\caption{Positive contrast bounds and complete-action gaps}\label{tab:date-bounds54}
{\small\setlength{\tabcolsep}{4pt}
\begin{tabular}{rrrrrr}
\hline
$T$ & Method & Pass & Date & $\chi$ upper & $\varepsilon$ upper \\
\hline
2 & W & 1 & 0 & 0.82795 & 0.18860 \\
2 & W & 1 & 1 & 0.64519 & 0.02510 \\
2 & W & 2 & 0 & 0.78262 & 0.05797 \\
2 & W & 2 & 1 & 0.65789 & 0.01568 \\
2 & F & 1 & 0 & 0.78145 & 0.05427 \\
2 & F & 1 & 1 & 0.65671 & 0.01605 \\
2 & F & 2 & 0 & 0.78129 & 0.04813 \\
2 & F & 2 & 1 & 0.65789 & 0.01564 \\
3 & W & 1 & 0 & 0.88888 & 0.28837 \\
3 & W & 1 & 1 & 0.80711 & 0.11312 \\
3 & W & 1 & 2 & 0.65789 & 0.02092 \\
3 & W & 2 & 0 & 0.86841 & 0.22498 \\
3 & W & 2 & 1 & 0.78262 & 0.04873 \\
3 & W & 2 & 2 & 0.65789 & 0.01568 \\
3 & W & 3 & 0 & 0.82423 & 0.16149 \\
3 & W & 3 & 1 & 0.78262 & 0.04873 \\
3 & W & 3 & 2 & 0.65789 & 0.01568 \\
3 & F & 1 & 0 & 0.82423 & 0.04903 \\
3 & F & 1 & 1 & 0.78262 & 0.03318 \\
3 & F & 1 & 2 & 0.65789 & 0.01521 \\
3 & F & 2 & 0 & 0.82423 & 0.04903 \\
3 & F & 2 & 1 & 0.78262 & 0.03318 \\
3 & F & 2 & 2 & 0.65789 & 0.01521 \\
3 & F & 3 & 0 & 0.82423 & 0.04903 \\
3 & F & 3 & 1 & 0.78262 & 0.03318 \\
3 & F & 3 & 2 & 0.65789 & 0.01521 \\
\hline
\end{tabular}}
\par\smallskip
{\footnotesize The symmetric bound is obtained from anchored signed endpoints and the independent scalar fallback. The deployed gate retains the signed endpoints. The directed gap already includes whole-cell evaluation, continuous-action covering, integration and arithmetic effects; these are not added a second time. Per-bin and per-candidate endpoints remain in the compressed raw records.}
\end{table}

\begin{table}[htbp]
\centering
\caption{Directed, symmetric and scalar candidate-gate comparison}\label{tab:gates54}
{\small\setlength{\tabcolsep}{4pt}
\begin{tabular}{rrrrrrr}
\hline
$T$ & Method & Pass & Directed & Symmetric & Scalar & Blocked \\
\hline
2 & W & 1 & 1657 & 525 & 0 & 15857 \\
2 & W & 2 & 765 & 63 & 0 & 16254 \\
2 & F & 1 & 656 & 108 & 0 & 17722 \\
2 & F & 2 & 0 & 0 & 0 & 17725 \\
3 & W & 1 & 1208 & 440 & 0 & 25243 \\
3 & W & 2 & 454 & 56 & 0 & 25464 \\
3 & W & 3 & 203 & 0 & 0 & 25446 \\
3 & F & 1 & 64 & 64 & 0 & 25778 \\
3 & F & 2 & 0 & 0 & 0 & 25778 \\
3 & F & 3 & 0 & 0 & 0 & 25778 \\
\hline
\end{tabular}}
\par\smallskip
{\footnotesize Counts concern non-incumbent candidate comparisons. Directed counts use strict negative upper advantages. Symmetric and scalar counts use their recorded weak gates. Blocked counts use the directed rule. All three are evaluated on the same proposed actions; their counts are not independent policies.}
\end{table}


\subsection{What the verification gauge computes}
The choice $h=0$ makes the residual at an intermediate date equal to the incumbent stage cost and the terminal residual equal to $g$. Thus a pairwise continuation calculation is a genuine evaluation of policy-value differences from the primitives. Its residual-difference enclosures are generally nonconstant and its continuation-contrast allowances generally positive. A signed interval for the advantage is more informative than replacing that interval by its absolute maximum before deciding whether an action helps. The symmetric and scalar candidate counts in the table are computed on exactly the same menus and whole cells.

The inherited code also records the largest absolute stage-residual and propagated pair endpoints visited by the recursion. These are diagnostics of the visited boxes, not a claim that a single sampled number is the globally defined function $H_t$ or $D_t$. The certified functions are supplied by the enclosure procedure on any requested state-pair box and are used on all boxes needed by the exhaustive root-cell and action-cover calculation. Neither smoothness nor interpolation of the discontinuous actor is assumed. Future rectangle extrema include every intersected acquired cell, with both closed endpoints included.

\subsection{Coupling, integration, and arithmetic}
Each dyadic innovation bin has its exact probability under the original continuous law. Both policy trajectories use the same innovation inside that bin; the joint law is a coupling and each marginal remains the stated economic law. Interval evaluation encloses the bin's entire contribution before its probability is applied. No bin midpoint is treated as an exact conditional expectation. The last innovation uses the analytic final-stage formula in the original two-state model.

Centered transition and cost identities reduce interval dependency before numerical rounding. A separately propagated enclosure of the state difference is intersected with other valid information. The retained tests compare primitive expressions with rational identities, check actor rectangles against brute-force extrema, verify the continuous final-cost formula and exercise the full finite-horizon directed recursion in exact finite models. Source hashes bind these algorithms to the deposited records. Replaying stored endpoints is not an independent derivation of every endpoint from primitives; the release makes that distinction explicit.

\section{Complete work and storage records}\label{supp:work54}
\begin{table}[htbp]
\centering
\caption{Executed construction, verification work and storage}\label{tab:work54}
{\small\setlength{\tabcolsep}{4pt}
\begin{tabular}{rrr r r r r r}
\hline
Cohort & $T$ & Method & Build (s) & Total (s) & Pair nodes & Table bytes & RSS (KiB) \\
\hline
P & 2 & W & 0.260 & 6.612 & 304363 & 299520 & 41136 \\
P & 2 & F & 0.271 & 6.643 & 312605 & 299520 & 39632 \\
P & 3 & W & 1.975 & 149.587 & 4931489 & 449280 & 50436 \\
P & 3 & F & 1.583 & 149.178 & 4886034 & 449280 & 44716 \\
A & 2 & W & 0.252 & 6.646 & 304363 & 299520 & 41124 \\
A & 2 & F & 0.256 & 6.579 & 312605 & 299520 & 40372 \\
A & 2 & A & 12.346 & 18.521 & 301893 & 299520 & 41356 \\
\hline
\end{tabular}}
\par\smallskip
{\footnotesize P denotes the primary cohort and A the separately executed adaptive cohort. Total is the recorded service clock through its final durable record, excluding shared inference. Pair nodes refer to the last pass, not the whole service. Table bytes are logical integer-table storage; RSS is the process peak. Neither quantity is a complete hardware-independent complexity measure. CPU affinity and library versions are in each service record.}
\end{table}

The service record distinguishes primitive construction, acquired-actor preparation, each pass, serialized endpoint arrays and the final durable record. The shared path-cost service has its own clock, policy identities, sample stream and endpoint archive. The catalogue table in the article conservatively charges every own executed pass and the full shared inference service. Postprocessing, replay, typography checks and clean-archive reproduction have separate measured clocks in the publication record; they are not silently assigned zero cost.

The logical actor-table count differs from total process memory. At five bits the exact two-dimensional minimum and maximum sparse tables require 149,760 bytes per actor, including its action array and query-log array. The process also stores interval batches, policies, retained outputs, interpreter state and numerical-library workspaces. Its reported peak resident set therefore exceeds the actor-table count. The pair-node count in the table is the final pass count; every earlier pass has a separate raw work record and must be included when summing total verification operations.

Adaptive refinement is paid for at every attempted construction rung, including pilots. The surplus-driven comparator's recorded date models are nonuniform and its full adaptation history remains in its service JSON. A tensor layout with unequal coordinate spacings is still a tensor layout. Neither this storage account nor the common-innovation recursion is evidence for a non-tensor high-dimensional runtime claim.

\section{Approximately null coefficients and controlled errors}\label{supp:nonnull54}
Let the perturbed action exposures be
\[
 b_1=\frac12+\xi+\lambda,\qquad
 b_2=\frac14-\frac{\xi}{2}-\lambda,
\]
and the ReLU direction be $v=(1,-2-\eta)$. The exact action coefficient is
\[
 s=b_1-(2+\eta)b_2
   =2\xi+3\lambda-\frac{\eta}{4}
       +\frac{\eta\xi}{2}+\eta\lambda.
\]
For a common innovation, the one-Lipschitz property gives the verified continuation allowance $|Ms|\,|a-b|$. The gate adds its discounted upper enclosure. Zero is used only after the coefficient is checked exactly. The parameter $\lambda$ is a known action-dependent shift of the innovation mean, represented in the actual perturbed kernel; it is not a test that an arbitrary unknown coupling has the correct marginals.

For each of the 384 cases, the raw archive retains the true-cost interval, nuisance interval, allowance, critic upper bound and corrected upper bound for every candidate. It also retains both the safe selected policy and the deliberately uncorrected false-null policy. The replay reconstructs the choices, verifies their identities and counts strictly harmful false-null actions only when their true-cost lower endpoint is positive. The primary economic experiment remains at the original unperturbed law. Repetition of one cell under different amplitudes or perturbations is not a new independent policy-cost observation.

\section{Cost inference, selection, and release scope}\label{supp:inference54}
All marginal cost and paired difference intervals are recomputed from the deposited outward path endpoints using the frozen bounded-sample formula. Raw path gains are allowed to be negative: conditional-expectation policy improvement does not imply pathwise dominance under a common shock. The support used for a paired gain remains the signed support $[-H,H]$, rather than clipping away adverse paths. Identical policy arrays have exact zero contrasts and are recorded as identities.

Selecting a pass with the smallest upper endpoint is valid on the simultaneous event covering the frozen finite catalogue. It does not permit new policies, undeclared initial laws or additional estimands without a corresponding coverage account. The two R53 statistical families remain separate, and the R54 replay adds no sample size. The pseudorandom seed fixes reproducibility; independence is the statistical model, not a theorem about the seed.

The release preserves every inherited labeled result in the active and complete documents and verifies unchanged bytes for historical evidence and frozen scientific sources. It compiles the main article, supplement, response and both complete editions, checks references and typography, and rebuilds from a clean archive without scientific reruns. Compilation, finite tests and source identities are not an external mathematical acceptance decision. The original broader NBO applications retain their own hypotheses; the new linear action-kernel identities concern the expected-cost model actually stated in the article.

\bibliographystyle{ecta-fullname}\bibliography{references}
\end{document}
